\documentclass[11pt,a4paper]{article}

\usepackage[english]{babel}
\usepackage[T1]{fontenc}
\usepackage{lmodern}
\usepackage[a4paper,margin=1in]{geometry}

\usepackage{xcolor}
\usepackage{graphicx}
\usepackage{float}
\usepackage{booktabs}
\usepackage{array}
\usepackage{caption}
\usepackage{subcaption}
\usepackage{tabularx}

\usepackage{amsmath,amssymb,bm}
\usepackage{enumitem}
\usepackage{titlesec}
\usepackage[authoryear]{natbib}
\usepackage{hyperref}
\usepackage{fancyhdr}
\usepackage{pdflscape}
\usepackage{xfrac}

\usepackage{tikz}
\usetikzlibrary{arrows.meta, positioning, shapes.geometric}

\newtheorem{assumption}{Assumption}

\definecolor{darkblue}{rgb}{0,0,0.8}
\hypersetup{
	colorlinks=true,
	citecolor=darkblue,
	linkcolor=darkblue,
	urlcolor=darkblue
}

\setlength{\parskip}{0.55em}
\setlength{\parindent}{0pt}
\begin{document}


\tableofcontents

\pagebreak

%<*ch3body>

% ============================================================
% EDITORIAL PASS v1 — 2026-04-09
% Agent: Ch3 editor
% Source: Cosmic Letter (Ch3_Cosmic_Letter_FromCh2.md)
% Scope: Introduction (new), S1–S3 (modified), S4 (new, stacked)
% ============================================================

% ============================================================
% INTRODUCTION — 100% NEW
% Collects three K5 mandates from Ch2.
% Positions leverage ratio as the answer to institutional causality.
% Frames Chile as the structural test case, not the extension.
% Places the three-chapter over-accumulation spine.
% ============================================================

\section{Introduction}
\label{sec:introduction}

% [K5 MANDATE: US Stage B closer — direction of institutional causality]
% [K5 MANDATE: Conclusion para 2 — causal architecture of Fordist settlement]
% [K5 MANDATE: Chile Stage B closer — BoP constraint converting technological dependency into financial fragility]
% COSMIC LETTER: Discovery 2 — three-node over-accumulation spine

Chapter~2 established the anatomy of Fordist exhaustion and the structural
conditions of peripheral over-accumulation. The four-channel Weisskopf
decomposition made legible the compositional structure of profitability
decline at both the centre and the periphery, and the threshold vector
error-correction model identified the balance-of-payments constraint as a
binding regime in Chile's import--accumulation dynamics. What Chapter~2
could not establish---and what no accounting decomposition can
establish---is the monetary architecture through which the
balance-of-payments constraint reproduces itself as a law of motion: the
leverage ratio through which world money enters peripheral accumulation,
converts external balance into a state variable, and transforms the
technological trap identified in Chapter~2 into the financial
subordination that is the structural condition of peripheral capitalism.

Three objects were passed forward explicitly. First, the US Stage~B
closer named the direction of institutional causality as structurally
open: whether demand activation counteracted over-accumulation, or
whether over-accumulation made demand activation structurally necessary,
is a question no profitability accounting can resolve. Second, the
Conclusion recognised that the four-channel decomposition makes the
anatomy legible but that the causal architecture of the Fordist
settlement---and its peripheral counterpart---requires a framework
beyond the decomposition itself. Third, the Chile Stage~B closer
identified the structural dynamic through which the balance-of-payments
constraint converted technological dependency into financial
fragility---the path-dependent escalation from the mechanisation trap to
the debt trap that the decomposition exposes but does not formalise. Each
of these objects finds its resolution in the same theoretical
construction: the leverage ratio as the law of motion of peripheral
over-accumulation in monetary form.

% COSMIC LETTER: Discovery 2 — over-accumulation spine: theta < 1 -> theta(omega) -> LevR
The dissertation's analytical spine runs through three nodes. Chapter~1
established that $\hat{\theta} < 1$---the transformation elasticity as a
qualitative invariant of capitalist technical change, the scalar
condition under which mechanisation necessarily compresses the profit
rate. Chapter~2 gave the invariant its distributive conditioning:
$\hat{\theta}(\omega) = 0.918$ for the United States, and
$\hat{\theta}^{CL}(\omega, \varphi) < 1$ for Chile under the
balance-of-payments trap---the same scalar, now distribution-conditioned
and BoP-constrained. This chapter completes the spine. The leverage ratio
is the monetary form of what $\hat{\theta}^{CL} < 1$ has been all along:
the balance-of-payments constraint makes over-accumulation
self-reinforcing at the material level; the leverage ratio formalises how
that self-reinforcement reproduces itself in the monetary register---through
reserves, exchange rate, external debt service, and the domestic monetary
expansion that the distributive conflict itself generates.

% COSMIC LETTER: Discovery 1 — Chile as structural test case, not extension
Chile is not an extension of the US case. It is the structural test
case. Chapter~2's methodological redesign---replacing the US-only ARDL
investment function with a symmetric causal inference strategy applicable
to both cases---produced a consequence the redesign itself did not
name: Chile's balance-of-payments-constrained structural position now
generates the necessary-condition test that the US case alone cannot
provide. The case where the BoP constraint is operative lets us identify
what the constraint does when absent. The US Fordist settlement
operated within a demand corridor unconstrained by external balance; the
Chilean case operated within a corridor whose width was set by the
leverage ratio. Comparing the two is not extending the analysis---it is
completing it.

The chapter proceeds in four sections. Section~\ref{sec:A_contested}
reviews four schools of economic thought on the Unidad Popular---classical
dependency theory, the neoliberal school, neo-structuralism, and
Historical Political Economy---through a diagnostic that traces what each
school cannot ask and why, closing the circle with the revival of
dependency theory as a research programme equipped with formal tools.
Section~\ref{sec:B_FS_WM} establishes the monetary architecture:
the commodity-to-credit transition, the dollar as world money under the
monetary expression of value, the Bretton Woods institutionalisation of
peripheral subordination, and financial subordination as a state
variable. Section~\ref{sec:C_levr} formalises the leverage ratio---its
constitution as a balance sheet object, the law of motion governed by
the balance-of-payments identity and the monetary expansion multiplier,
and the bridge from theory to identification through the two state
variables that condition the transmission channel.
Section~\ref{sec:empirical_strategy} operationalises the formal apparatus
as an empirical strategy: state-dependent local projections organised
around three links---external transmission, domestic propagation, and
regime-conditioned response---anchored in Robinson's Historical Time
principle that causal compounding requires temporal co-activation within
the same calendar-year window.


% ============================================================
% SECTION 1 — CONTESTED APPROACHES
% Modifications: S1.4 (s_pi insertion), S1 close (transitional paragraph)
% All other subsections UNCHANGED.
% ============================================================

\section{Contested Approaches to the Political Economy of the Rise and Fall of the Unidad Popular}
\label{sec:A_contested}

Four schools of economic thought have produced accounts of the rise and fall of Chile's Unidad Popular government. Each emerged under specific political-economic conditions that shaped not only its conclusions but the questions it could ask. The neoliberal school crystallised within the networks that linked the Pontificia Universidad Católica de Chile to the University of Chicago and gained state power through the dictatorship that followed the 1973 coup. The neo-structuralist school emerged in the CIEPLAN research centres during that same dictatorship, positioning former structuralists as the intellectual core of the post-dictatorship Concertación governments. Historical Political Economy consolidated institutionally with the credibility revolution's application of causal identification methods to historical questions. And dependency theory --- the tradition that produced the intellectual foundations of the UP programme itself --- was physically destroyed by the coup before resolving the theoretical limitations its own practitioners had identified.

What matters is not which school is right in isolation but what each one cannot ask --- and why. This section reviews these four schools through a heuristic matrix cross-referencing each with eight analytical dimensions: historical context of elaboration, main hypothesis, summary of the political economy of the rise and fall, key references, theory of aggregate demand and distributive conflict, strengths, limitations, and overall assessment. The matrix appears at the section's close as a visual anchor. But the narrative does not follow the matrix's grid. It follows an arc.

The arc is circular. It opens from within the dependency tradition --- the intellectual world the UP programme emerged from --- traces that tradition's political and intellectual defeat, examines the accounts that filled the vacuum or evaded the questions the defeat left unanswered, and returns to the dependency research programme renewed with formal tools the classical version lacked. The recurring absence across all four schools is class conflict as endogenous to macroeconomic dynamics. Each school fails to integrate it for structurally different reasons. The circle closes when the renewed programme addresses this absence --- connecting external transmission mechanisms, internal distributive conflict, and macroeconomic performance as interacting dynamics within a single framework. What follows is that circle.


\subsection{Classical Dependency Theory: ``The dependency school and its defeat''}
\label{sec:A1_classical_deptheory}

% [UNCHANGED — see original, lines 66–99]
% Full content preserved from Chapter3_PEUP.tex lines 66–99.

\subsubsection{Seminal Contributions to Classical Dependency Theory }

Dependency theory did not comment on the Unidad Popular from a distance --- it built the intellectual world from which the UP programme emerged. The tradition's foundations run through two lineages that converged, institutionally and theoretically, in Santiago during the long sixties. The Monthly Review school --- \citet{Baran1957}'s account of the economic surplus and its systematic appropriation under peripheral capitalism --- provided the underconsumptionist framework that \citet{Frank1967} radicalised into the development-of-underdevelopment thesis: peripheral underdevelopment is not a stage on the road to development but an active product of integration into the world capitalist system. Simultaneously, Prebisch and the ECLAC structuralists established the centre-periphery framework and the deteriorating terms-of-trade hypothesis that would become the terrain on which dependency theory was constructed --- \citet{Prebisch1981}'s mature work already in dialogue with the dependency critique he had helped provoke.

The synthesis of these lineages happened at the Centro de Estudios Socioeconómicos (CESO) at the Universidad de Chile, directed by Theotônio dos Santos. The CESO network --- Frank, \citet{DosSantos1970}, \citet{Bambirra1971}, \citet{CaputoPizarro1970}, with Marini working in close dialogue --- elaborated the theoretical core: a structurally subordinated bourgeoisie incapable of leading autonomous development, imperialism and external constraint as constitutive of peripheral economic dynamics, and socialist transformation as the only viable developmental road. Bambirra's typology of dependent capitalism positioned Chile as a Type I case --- early primary-export integration --- giving the abstract dependency hypothesis its comparative empirical form. Dos Santos's concept of a \textit{``new dependence''} specified the financial-industrial subordination mechanism. Caputo and Pizarro's \textit{Imperialismo, dependencia y relaciones económicas internacionales}, produced directly within the CESO, grounded the framework in a systematic account of international economic relations from classical Marxism --- the intellectual scaffolding for the UP's banking nationalization strategy. These were not academic exercises. They were the theoretical foundations of a programme of government, produced inside the institutional spaces that programme created \citep{Kay2010, CardenasLanaSeabra2022}.

The tradition's internal tensions were already active before the coup tested them. \citet{Marini1973}'s \textit{Dialéctica de la dependencia} pushed the analysis toward superexploitation: peripheral capital compensates for unfavourable terms of trade by depressing wages below the value of labour power, producing a dependent reproduction schema in which internal consumption is systematically compressed. This was the tradition's most developed account of distributional conflict --- and the one most directly relevant to the UP's wage dynamics. But Marini's world-market overdeterminism --- deriving domestic class dynamics from the global logic of capital --- sat in tension with \citet{CardosoFaletto1979}'s structural-historical method, which insisted that dependency takes different concrete forms depending on the specific class configuration within each social formation, and that the state and domestic political alliances are analytically autonomous. The disagreement was not peripheral. It concerned whether domestic political mediation was an independent object of analysis or a derivative of world-market imperatives --- a question the tradition raised with precision but never settled.

\citet{Bambirra1974}'s \textit{La revolución cubana: Una reinterpretación}, written during her years in Chile, made the stakes of this unresolved tension explicit by introducing a strategic mirror. Against what reference point was the UP's peaceful road being theorised? The Cuban Revolution was the implicit answer, and Bambirra's text stated it directly --- dependency theory in practice, producing comparative analysis of revolutionary strategy from within one of the processes it sought to explain. These were the CESO network's intellectual interventions \textit{during the UP itself}: not retrospective scholarship but real-time theoretical production embedded in the political conjuncture.

On the margins of the CESO tradition but feeding into its conceptual vocabulary, a parallel cluster theorised the structural labour-market substrate on which the UP's social base rested. \citet{Nun1971}'s concept of the relative surplus population --- the marginal mass that peripheral capitalism cannot absorb into formal employment, unlike the cyclical reserve army in classical Marxism --- identified the structural condition of the urban poor from whom the UP coalition drew much of its political support. \citet{Quijano1974} mapped the marginal pole of the dependent economy; Castells applied the framework to Santiago's \textit{pobladores} movement as a class actor in real time; Portes specified the informal sector's articulation with the world economy through subsidised labour; Tokman modelled the intersectoral surplus transfer from informal to formal through a Sraffian lens. \citet{Kay1980}'s account of agrarian change and migration traced the rural dispossession that generated this urban relative surplus population. This cluster is not peripheral to the dependency argument --- it is the dependency argument applied to the labour market and the city.

\subsubsection{Classical Dependency Theory on the Political Economy of the Rise and Fall of the Unidad Popular}

The empirical literature on the UP itself is dense and multi-layered. \citet{DeVylder1976}'s \textit{Allende's Chile} remains the anchor reference: the most comprehensive political economy account of the UP's internal contradictions and external pressures from within the dependency tradition, written by a participant-observer economist. The expropriation strategy's internal contradictions --- unauthorised factory occupations outpacing programmatic control, the tension between socialising monopoly firms that were less labour-intensive and the small enterprises workers actually occupied --- and the U.S. credit blockade and spare-parts embargo are documented with primary data. De Vylder established the analytical template that would define the tradition's empirical contribution: internal class dynamics and external constraint as interacting forces, not competing explanations.

Around De Vylder, a cluster of real-time studies built the empirical panorama from different entry points. Stallings and Zimbalist provided the longitudinal comparative frame: distributional conflict across the Alessandri, Frei, and Allende governments as a continuous structural variable, not a UP-specific pathology \citep{StallingsZimbalist1975, Stallings1978}. \citet{EspinosaZimbalist1978} brought the analysis to the shop floor --- firm-level surveys showing worker participation improving productivity in the short run before macroeconomic deterioration set structural limits. \citet{Winn1986}'s micro-history of the Yarur textile workers' takeover demonstrated that labour radicalization was endogenous: workers pushed the UP programme beyond its designed boundaries from below, not through agitation from above. \citet{Kay1978} traced the same dynamic in the countryside --- agrarian reform as class struggle, peasant mobilization driving expropriation beyond the programme's designed pace and scope.

The external constraint was documented with equal density. \citet{PetrasMorley1975} provided the most comprehensive empirical treatment of U.S. destabilisation --- CIA operations, ITT corporate intervention, multilateral credit blockade, military aid to the opposition --- tracing the relationship between the North American \textit{``imperial state''} and multinational corporations that resulted in the 1973 overthrow. \citet[ch.~5]{GriffithJones1981} specified the financial and monetary channel through which external constraint became internal crisis: the state-sector deficit, monetary expansion, and balance-of-payments deterioration as interacting variables, not sequential policy errors. \citet{Sideri1979}'s edited volume, produced from a seminar at The Hague gathering economists who had served in the UP government, is a unique source of insider self-criticism --- chapter by chapter covering the external sector, copper nationalisation, the industrial \textit{Área de Propiedad Social}, banking nationalisation, and inflationary dynamics.

The political debate within the Marxist left was itself documented in real time. \citet{SweezyMagdoff1974}'s collection assembled contributions by Sweezy, Petras, Zimbalist, and Stallings as the UP unfolded, including the Sweezy-Zimbalist exchange on whether the bourgeois state could serve as an instrument of socialist transition --- a dispute that named the tradition's state-theory problem in the language of strategic urgency. \citet{Zavaleta1974} theorised the conjuncture analysis by looking at the different trajectories of revolutionary movements in Chile and Bolivia in comparative and relational perspective, tracing the contradictory development of dual power --- the coexistence of two competing authority structures, the bourgeois state and the emerging popular institutions --- as the structural moment the UP could neither absorb nor resolve.

The retrospective synthetic treatments close the corpus. \citet{Kay2010} provides the most comprehensive overview of the tradition --- the reformist-Marxist debate within dependency theory, the critique that followed, and the UP as a unique case of a \textit{dependentista} government in practice. \citet{Palma1978} conducts the tradition's most rigorous internal critique, tracing the theoretical debate back to Bolshevik disputes and the establishment of Communist parties around a stageist political strategy anchored in the impossibility of development \textit{à la} Frank. His conclusion --- that dependency theory should be treated not as a closed formal theory but as a methodology for the concrete analysis of historical junctures --- is the single most productive verdict the tradition produced about itself, and the one this essay most directly inherits.

The tradition's strengths were inseparable from its conditions of production. It placed imperialism, external constraint, and structural bottlenecks at the centre of the analysis at a moment when these were not rhetorical categories but lived constraints: the copper terms of trade, the spare-parts embargo imposed by the United States, the financial blockade, the agrarian bottleneck \citep{DeVylder1976}. It motivated political economist and marxist around the world to produce real-time interventions and empirical work by participants and observers of the process itself --- a richness of sources unmatched by any other school reviewed in this chapter. And it insisted --- against what every subsequent school would elide --- that distributive conflict was endogenous to the economic dynamics of peripheral capitalism, while rejecting its portrait merely as a policy variable to be managed. The tradition had, in its classical form, no formal mechanism connecting external shocks to domestic distribution --- Frank's underconsumptionism and Marini's superexploitation thesis provided fragments but no integrated model --- yet its insistence on asking the question was itself a contribution that none of subsequent accounts from other theoretical traditions would match. One finding is uncontested across all four schools reviewed in this chapter: the UP's strategic incoherence: the gap between programme design and political implementation, and the lack of cohesion within the left at critical junctures while class struggle operational field intensified, pushing further the strategical disadvantage \textit{vis-a-vis} the recomposition of reactionary forces and Chile's ruling classes. What remains empirically unresolved is the causal weighting between financial-external constraint mechanisms and domestic policy choices --- a question the classical tradition raised but could not formally operationalized, nor subject to empirical scrutiny, which remain yet to be resolved.

Classical Dependency Theory suffered a double defeat --- political and intellectual. The coup of September 1973 physically destroyed the networks that produced dependency scholarship in Chile. The CESO was dismantled, its members expelled or imprisoned, its intellectual production scattered into exile. \citet{Marini1976}'s \textit{El reformismo y la contrarrevolución}, written from exile, is the primary testimony of this moment: the tradition reflecting on its own defeat from outside the country it had tried to transform --- reformism diagnosed as a class strategy, the counterrevolution as restoring the conditions of dependent accumulation. The political destruction did not come from intellectual inadequacy. But it coincided with unresolved theoretical limitations that the tradition's enemies did not need to expose --- they were already visible from within.

There was no theoretical unity across the tradition's three strands. \citet{Palma1978} had mapped the incompatibility between Frank's stagnationist thesis, the Cardoso-Faletto structural-historical method, and the ECLAC reformist current, and no synthesis followed. There was no relational theory of the state: the tradition oscillated between liberal structuralist treatments that located the state as an autonomous modernising agent and instrumentalist Marxist treatments that reduced it to a tool of the dominant class. Cardoso and Faletto's relational method --- insisting on the analytical autonomy of domestic political alliances --- became highly influential in academic circles in the north, but only revisited by Latin American scholarship from exile or during post-dictatorial period in marginalized circles of policy making. \citet{Zavaleta1974}'s concept of dual power named the gap with precision, identifying the coexistence of bourgeois state and popular institutions as the structural moment the UP confronted, but from outside the tradition's theoretical core. And the overdeterminist world-market schema --- Marini's superexploitation deriving domestic class dynamics from the global logic of capital --- left domestic political mediation underdetermined, a problem the Sweezy-Zimbalist exchange exposed in real time as a dispute over the state's class character \citep{SweezyMagdoff1974}.

The dependency programme was not refuted. It was interrupted --- politically, by the destruction of its institutional base, and intellectually, by the unresolved questions it had raised but not yet answered. But the programme's enemies did not wait for answers. While the tradition was still assembling its theoretical foundations, the Chicago Boys had already written theirs --- not as a retrospective but as a prospectus, ready before they had state power. The victors did not write history after they won. They wrote the programme before they seized the state.


\subsection{The Rise of the Chicago Boys and the Neoliberal School}
\label{sec:A2_neoliberal_chicago}

% [UNCHANGED — see original, lines 102–124]

The neoliberal school did not emerge from within the intellectual debate on the Unidad Popular --- it emerged from its political defeat. The Chicago Boys' trajectory from academic formation to state power is the best-documented case of transnational ideological transfer in Latin American economic history. The Pontificia Universidad Católica de Chile signed an exchange agreement with the University of Chicago's Department of Economics in 1956 --- a Cold War instrument, funded initially through the International Cooperation Administration and later through the Ford Foundation, designed to train a cadre of economists in monetarist orthodoxy who would displace the structuralist consensus from within Chile's own academic institutions \citep{Fischer2009, Valdes1995}. By the late 1960s, the programme had produced a network of PUC-trained economists fluent in Friedman's monetarism and Harberger's public finance, professionally connected to each other through graduate cohorts in Chicago, and politically connected to the Chilean right through Jaime Guzmán's gremialista movement. The crisis of right-wing electoral politics across the 1960s --- from Alessandri's failed stabilisation to the Christian Democrats' agrarian reform --- radicalised this network toward a comprehensive programme for the restructuring of Chilean capitalism.



That programme was \textit{El Ladrillo}. \citet{DeCastro1992}'s edited volume was not a retrospective analysis of the UP --- it was written as the economic programme for what to do after the UP was overthrown. The title is revealing: ``The Brick'' names itself as a foundation, not a commentary. Its diagnostic identified six interrelated pathologies of the developmental state: low growth from decades of failed state-led programmes, suffocating statism crowding out private activity, productive employment scarcity from market distortions, conflict-driven inflation from the over-politicisation of economic life, agricultural stagnation burdening the balance of payments, and persistent poverty from ineffective social policies. The core hypothesis --- a secular tendency of the Chilean developmental state toward socialisation of the economy --- framed the UP not as a rupture but as the culmination of a long-running process. Under this reading, Allende's government merely accelerated what Alessandri and Frei had already set in motion.

The analytical structure deserves recognition for what it was: a counterrevolutionary vanguard text that took macroeconomic magnitudes seriously. \citet{Clark2017} traces the institutional trajectory from the PUC-Chicago network through the planning committees that assembled \textit{El Ladrillo} during the UP period itself, showing that the document circulated among military plotters before the coup --- the economic programme and the political intervention were coordinated, not sequential. This is not incidental context. The conditions of intellectual production determine what questions a school can ask and which it must suppress. A programme designed to justify the restructuring of an economy after a military overthrow cannot theorise the class conflict that made the overthrow necessary --- to do so would name the social forces its own sponsors represent. The omission is structural, not accidental.

The core hypothesis is empirically weak. The evidence for a secular tendency toward socialization across Chile's developmental period does not survive sustained scrutiny --- the trajectory from the 1930s through the 1960s shows fluctuation in state economic intervention, not monotonic expansion, and the comparative record across Latin America does not support the claim that Chile was exceptional in this regard. What \textit{does} survive is a set of auxiliary hypotheses. The argument that heavy price distortions throughout the developmental period favoured inefficient industry at the expense of agricultural development and export diversification corresponds to documented patterns in the data \citep{DeCastro1992}. The claim that exchange-rate overvaluation generated balance-of-payments stress is consistent with Chile's external accounts across multiple governments. These concessions are worth making explicitly --- and every empirically grounded one comes from the historically situated primary text, not from its theoretical apparatus. \citet{Espinosa2021}'s contemporary Austrian revival contributes zero independent empirical observations. The distinction matters: the school's empirical credibility resides in De Castro's diagnostic, not in the theoretical framework that claims to explain it. They demonstrate that the dependency tradition's defeat was not because the neoliberals had superior economic analysis --- some of their second-order observations were accurate --- but because they had state power. The core hypothesis provided the ideological architecture; the auxiliary hypotheses did the empirical work.

The school has no theory of distributive conflict. Distribution appears in \textit{El Ladrillo} as an outcome of market equilibrium --- get prices right, remove distortions, and the allocation of resources will reflect underlying scarcities. Demand management is either unnecessary, following Say's Law, or actively harmful, following Friedman's monetarism. The foundational premises --- that price distortions and state intervention cause economic crisis --- are treated as self-evident starting points rather than as claims requiring theoretical justification. Whether these premises derive from a particular moral vision of society or from economic theory proper is a question the literature on neoliberal economic thought has explored extensively, and the answer matters: if the premises are ultimately normative rather than scientific, then the diagnostic built upon them is a political programme presenting itself as economic analysis \citep{Espinosa2021}.

This is precisely what the contemporary Austrian revival makes visible. \citet{Espinosa2021} revisits the UP as a case study in the impossibility of socialism, drawing on Mises's calculation argument and Hayek's knowledge problem to argue that any attempt to replace the price mechanism with centralised allocation must fail. The argument conflates Hayek's contributions to economic theory --- the insight that dispersed knowledge is coordinated through market signals --- with his political philosophy, in which any departure from market coordination constitutes a step on the road to serfdom. Espinosa does not acknowledge this conflation, and the result is an argument that extrapolates from the UP's specific contradictions to a universal impossibility claim without engaging the historical conditions that made those contradictions explosive: the credit blockade, the spare-parts embargo, the CIA-backed destabilisation campaign, the October 1972 bosses' lockout. The UP is flattened into a laboratory demonstration of an axiom, not analysed as a historical conjuncture.

The systematic omissions follow from the school's theoretical commitments, not from carelessness. Class conflict is absent because the framework has no category for it --- distributional outcomes are price signals, not power relations. Infrastructure bottlenecks --- the agrarian supply constraint, the transport system's fragility exposed by the lockout --- are invisible because the model assumes flexible adjustment. U.S. intervention is unmentioned because acknowledging an external political actor capable of engineering economic crisis would undermine the premise that the crisis was endogenous to socialist policy. The omissions are therefore not gaps to be filled; they are constitutive of the school's explanatory logic. To include what is omitted would require abandoning the theoretical framework that produces the omissions.


The neoliberal school constitutes a revolutionary class project presenting itself as scientific inquiry. This verdict does not rest on the tradition's political affiliations --- every school reviewed in this chapter has political affiliations --- but on the relationship between its theoretical apparatus and the evidence it claims to explain. \citet{Clark2017} identifies the deeper symmetry: the Chicago Boys and the socialist revolutionaries of the 1960s and 1970s were united in their common desire to utilise the power of the state to carry out a revolutionary reconstruction of state and society --- the main divergence was whether the revolutionary road would be socialist or capitalist. At the overt level, the Chicago Boys' technical record was uneven --- the 1982 financial crisis forced the state takeover of the financial sector, socialising private debts to 40 per cent of GDP. At the covert level, the project succeeded: the construction of a hegemonic capitalist elite via privatisation as class formation, the ``seven modernisations'' as sociological transformation of identity formation \citep{Clark2017}. \citet{Bockman2019} adds a further dimension: the Chicago Boys narrative itself performs what Gibson-Graham called capitalocentrism --- by treating the neoliberal victory as inevitable and the socialist alternative as self-evidently impossible, the conventional origin story ``erases what was at stake'' in the transition period, actively obscuring the anti-authoritarian transnational socialisms that were historically available \citep{Bockman2019}. The analytical omissions are not gaps to be filled --- they are constitutive of the school's political function. The core hypothesis is not supported by the data. The auxiliary hypotheses that \textit{are} supported could be incorporated into any macroeconomic framework and do not require the neoliberal theoretical apparatus to sustain them. De Castro's text earns the qualified recognition of having taken the Chilean economy's structural imbalances seriously --- it was a counterrevolutionary vanguard document that diagnosed real problems in order to justify a specific political solution. Its intellectual descendants have not maintained even that standard.

If the neoliberal account is the victors' programme imposed as history, who challenged it from within the post-dictatorship order? The technopols.

\subsection{The Neo-Structuralist School: ``A progressive sanitisation of the surplus approach''}
\label{sec:A3_neo_structuralism}

% [UNCHANGED — see original, lines 127–152]

The neo-structuralist school is an intellectual accomplishment of the post-dictatorship order, and its relationship to the tradition it renovated is one of strategic sanitisation. The CIEPLAN research centres emerged under the dictatorship as a technocratic opposition --- former structuralists, many trained abroad during the exile years, who monitored the Chicago Boys' economic management and prepared an alternative programme for the transition to democracy. \citet{Silva1991} traces the institutional trajectory: influenced by a political context in which the language of class conflict was prohibited and by the graduate training many of their members received in exile, the Corporación de Investigaciones Económicas para Latinoamérica constituted a new hub of economic thinking capable of casting doubt on orthodox neoliberal policy while incorporating concerns about equity, social security, and poverty alleviation into a framework compatible with market-oriented development. The systematic crisis produced by orthodox neoliberal management in the early 1980s gave these \textit{``CIEPLAN monks''} their political opening. By the time the Concertación governments took office, the former opposition technocrats had become the intellectual core of the new ruling coalition --- technopols administering neoliberal capitalism under democratic cover \citep{Sklair2000,Leiva2008}. The CIEPLAN economists were not simply former structuralists who moderated their politics under duress --- they were the Chilean fraction of a technocratic class whose professional authority depended on converting distributional questions into governance problems \citep{Maillet2016}.

The theoretical operation that made this trajectory possible was therefore the abandonment of classical structuralism's surplus approach. The selective inheritance is visible in real time. \citet{RosensteinRodan1973}'s engagement with the Chilean conjuncture --- a talk delivered in November 1972, while the UP still governed, from Boston University's Center for Latin American Development Studies that he had founded after years of collaboration with ODEPLAN under Frei --- already reads the crisis through the demand-management lens that NST would later canonise: a Keynesian short-run programme that won the battle of consumption and lost the battle of production, with no theory of the external constraint, no account of U.S.\ destabilisation, and the explicit verdict that the programme resembled populism more than socialism. The surplus approach is absent not because it was excised but because it was never part of the frame. What NST formalised through \citet{DornbuschEdwards1991} was therefore not a retroactive reinterpretation. It was the codification of a reading already available within development economics --- one in which the external constraint had no political economy and distributional conflict was a policy error, not an endogenous dynamic. Where Prebisch and the ECLAC tradition had centred the reinvestment of economic surplus and the structural conditions of peripheral accumulation, neo-structuralism relocated the analytical framework around market failure, macroprudential policy, and pragmatic state intervention --- a conceptual and policy framework where economic growth, equity, and democracy could be presented as mutually reinforcing without naming the class forces that determined which of these goals would prevail when they conflicted. Ffrench-Davis and Fajnzylber led the theoretical renovation within ECLAC itself, leaving import substitution behind and repositioning the commission as a producer of what \citet{Leiva2008} calls ``technologies of hegemonic discourse'' --- frameworks promoting prudential macroeconomic management and pragmatic policy intervention rather than challenging transnational capital as the original ECLAC had done. But \citet{Leiva2008}'s contribution is double. His positive argument --- under-deployed in the literature that invokes him --- is that the surplus approach had provided a theory of how the deterioration of the terms of trade is not merely an external shock but the form through which social surplus generated in the periphery is appropriated by the centre. Without it, the external constraint collapses into a balance-of-payments identity --- legible to Thirlwall but stripped of its class content. NST's abandonment of the surplus approach was therefore not merely a theoretical evolution. It was the abandonment of the only framework in which the external constraint had a political economy, not just a macroeconomic specification. The renovation was not a betrayal of structuralism's empirical insights. It was a betrayal of its theoretical foundations --- the surplus approach and the class-conflict categories that had made structuralism dangerous to the ruling classes that consolidated the institutionalisation of the dictatorial heritage during the post-dictatorship period.

The school's account of the UP crystallize itself around two complementary texts. \citet{Bitar2013}'s contribution, originally published in 1979, can be considered the first neo-structuralist understanding of the UP's political economy. As Allende's Mining Minister, Bitar wrote from the position of a participant, and his analytical ambition was to overcome the dichotomy between a financial approach to the balance of payments and a historical-structural one. His concept of a ``viability area'' --- a shrinking space within which the UP programme could have been implemented without triggering macroeconomic collapse, progressively narrowed by the interaction of structural reforms, political opposition, and external pressure --- remains the school's most sophisticated analytical construction. The integration of monetary, structural, and political dimensions in a single narrative is a genuine contribution, and the empirical detail Bitar provides is superior to what most dependency accounts achieved on the same questions.

But Bitar's work also established a pattern that would define the school's limitations. Despite multiple editions of his original contribution --- the most recent in 2013 --- Bitar never once engaged with the dependency theory literature that had developed in the 1970s on exactly the same questions \citep{Kay2010}. \citet{DeVylder1976}'s \textit{Allende's Chile}, the most comprehensive political economy account of the UP from within the dependency tradition, is absent from Bitar's analytical engagement across four decades and multiple editions. This is not an oversight. It is therefore the foundational omission around which the entire neo-structuralist account of the UP is constructed: the dependency tradition's arguments were equally consistent with the same data Bitar analysed, but engaging them would have required naming the class-conflict dynamics that the post-dictatorship settlement needed to exclude from legitimate economic debate.

The omission became paradigmatic through \citet{LarranMeller1991}. Their four-phase periodisation of the UP --- successful beginning in 1971, emergent disequilibrium in late 1971, decline and collapse through 1972--73, and unavoidable military intervention given macroeconomic mismanagement --- absorbed the case into the \textit{``populist macro paradigm''} that \citet{DornbuschEdwards1991} were simultaneously constructing as a regional template to dismantle Prebisch's legacy across the region. The paradigm defined economic populism as an approach emphasising growth and income redistribution while disregarding the risks of inflation, fiscal deficits, external constraints, and the reaction of economic agents to non-market policies. Three initial conditions were specified: economic stagnation coupled with popular frustration over inequality; the existence of idle capacity and international reserves creating apparent room for expansion; and the articulation of a programme promising real wage increases and economic reactivation through redistribution. The cycle then unfolds predictably: short-run success, reserve depletion, inflationary spiral, collapse.

The analytical power of this construction should not be underestimated. The core hypothesis --- that the UP's populist and socialist character was the engine of the crisis --- fits the available data on macroeconomic aggregates during the 1970--73 period, and its achievement of hegemonic status within progressive political discourse is not accidental. The four-phase model provides a clean narrative with identifiable turning points, quantifiable indicators, and a policy lesson that reads forward: re-distributive macroeconomic policy is historically doomed to failure. This is precisely its political function. By defining the problem as macroeconomic populism, the paradigm forecloses the possibility of wage-led growth regimes by definition. Any programme that uses macroeconomic policy for re-distributive goals is classified as populist before its results are examined. The category is diagnostic and normative simultaneously --- it describes a pattern and condemns it in the same gesture.

\citet{Meller2000}'s subsequent work on the UP consolidated this framework as the canonical NST economic narrative. His account added political polarization as an explanatory variable --- the radicalisation of both left and right narrowing the space for pragmatic management --- while maintaining the macroeconomic populism framing as the structural explanation. The result was a two-pillar architecture for the post-dictatorship centre-left's relationship to its own history: Bitar's political centrism, which located the UP's failure in insufficient pragmatism, and Meller's macroprudential framework, which located it in the irresponsibility of redistributive ambition. Together, these pillars supported a reading of the UP that was simultaneously self-critical and exculpatory --- the programme failed because of its own excess, not because of the structural conditions the dependency tradition had identified or the external intervention the empirical record documents.

The school's limitations are not gaps --- they are \textit{strategic omissions}. Larraín and Meller reference some of the neo-Marxist literature, including De Vylder, but do not engage with its core arguments on class-conflict dynamics. The dependency tradition's account of the credit blockade, the spare-parts embargo, and the CIA-backed destabilisation campaign is not refuted --- it is bypassed. The structural bottlenecks that De Vylder had documented as interactions between internal contradictions and external pressures are flattened into policy errors within the populist macro model. The October 1972 bosses' lockout --- which the dependency tradition analysed as a decisive moment of class confrontation --- appears in the neo-structuralist accounts as evidence of political polarization rather than as an expression of distributional conflict with determinate economic causes and consequences. The question of wage-led growth --- whether the UP's re-distributive strategy could have been sustainable under different external conditions --- is ruled out by the paradigm's own definition of populism, not by empirical investigation. The displacement was not abstract. Two economists named Ramos occupy symmetric positions in the same debate: Joseph Ramos, asked why demand expansion overheated --- a question of policy calibration. Sergio Ramos, from CESO and the UP government itself, asked why structural transformation was blocked --- diagnosing the same inflation as the product of class struggle and insufficient revolutionary transformation, not policy excess. The neo-structruralist framework did not merely omit the dependency tradition's arguments; it replaced an existing structuralist analysis available from within the UP's own intellectual production with a depoliticised substitute.

The neo-structuralist school functions as the ideological apparatus for post-dictatorship neoliberal hegemony under a progressive umbrella \citep{Leiva2008}. This assessment does not rest on the school's political affiliations --- Bitar and Meller were genuine opponents of the dictatorship, and their intellectual contributions carry real analytical substance. It rests on the relationship between what the school includes and what it excludes.

The omissions are not accidental but forensic --- systematic, traceable, and functional. \citet{Silva1991} demonstrates that the technocratic form itself depoliticises, making CIEPLAN monks and Chicago Boys structural equivalents in institutional function. \citet{Maillet2016} quantify the trajectory: four phases from monastery to technopols, producing four Finance Ministers and four Central Bank presidents from a single research institute --- political strategy trumped academic ambition almost from the outset. \citet{Sklair2000} locates this trajectory within world-system class formation: Foxley, like Cardoso, is not merely a technopol but a member of the transnational capitalist class whose professional authority depends on managing peripheral insertion into global capital circuits rather than contesting them --- former anti-imperialists turned administrators of the order they once opposed. \citet{Bockman2019} provides the genealogical evidence that sharpens the verdict further: CEPLAN staff --- the same economists who reconstituted themselves as CIEPLAN after 1973 --- had studied Yugoslav worker self-management socialism, hosted Horvat, and published on decentralised socialism as an alternative to both Soviet planning and capitalist markets. The abandonment of the surplus approach was therefore not an inheritance from a tradition that lacked it but a deliberate excision by economists who had once held it.

The inclusion of monetary, structural, and political variables in an integrated narrative is a real achievement. The exclusion of class conflict as an endogenous economic dynamic --- not a political externality but a force that shapes profitability, investment, and accumulation from within --- is the operation that makes the school serviceable to the process of depolitisation of the post-dictatorship institutional landscape. The sanitisation of classical structuralism was not an intellectual failure. It was an intellectual accomplishment in the service of a political project that needed the UP to be a cautionary tale about populist excess, not a case study in the structural contradictions of peripheral capitalism under external constraint.

That is why the neo-structuralist narrative infrastructure did not remain confined to the progressive centre-left. It became the default macroeconomic framing inherited by the next generation of scholarship --- including the field that promised to bring rigorous methods to historical questions.

\subsection{Historical Political Economy --- ``Methods without foundations''}
\label{sec:A4_historical_PE}

% [UNCHANGED through the main body — see original, lines 154–168]

Historical Political Economy is a genuine intellectual achievement --- the application of the credibility revolution's causal identification methods to questions that economists had largely ceded to historians and political scientists. The field consolidated institutionally around the \textit{Journal of Historical Political Economy}, and the \textit{Oxford Handbook of Historical Political Economy} \citep{JenkinsRubin2024, CharnyshFinkelGehlbach2023}, establishing rigorous methodological standards for empirical work on historical inquiry. The ambition is substantial, and the methods represent a real advance in the rigour with which historical questions can be addressed empirically, both in terms of data collection and empirical identification.

The bibliographic record is instructive. Badia-Miró and Llorca-Jaña have reconstructed long-run price series and living standards across Chile's developmental period; Arroyo Abad has traced wage inequality and labour-market dynamics across Latin America's long nineteenth century --- quantitative foundations that no other school reviewed in this chapter has matched in archival depth. On the causal identification frontier, \citet{AldunateGonzalezPrem2022} isolate the firm-level credit contraction produced by U.S.\ banking covert action during the Allende government through a difference-in-differences design; \citet{JaimovichToledo2021} trace the persistence of Mapuche territorial conflict to agrarian reform intensity during 1962--1973 using geocoded plot-level data and an instrumental variable strategy; \citet{GirardiBowles2018} identify the asset-price response to political shocks on the Santiago stock exchange through an event study. Rodriguez-Weber's distributional series establishes that Chilean inequality follows long waves driven by globalisation, institutional change, and coercive distribution --- a finding that moves the analysis beyond market mechanisms to state power as a determinant of distributional outcomes. The \textit{Journal of Historical Political Economy} programme's work on extractive institutions and institutional persistence extends this empirical ambition to the comparative Latin American frame. The field has therefore generated more empirical evidence on Chile's long-run political economy than any other school reviewed in this chapter --- and the theoretical framework organising that evidence is the neoclassical synthesis its parent disciplines provide. That is the constitutive boundary, not a gap.

The field's application to Chile reflects both its strengths and its structural limitation. \citet{RodriguezWeber2018}'s long-run distributional series is the strongest fit between HPE methods and the questions this essay poses. The work reconstructs 160 years of annual distributional data through two dynamic social tables --- 49 categories for 1860--1929, 116 categories for 1929--1970 --- supplemented by household surveys for 1957--2009, producing annual estimates of both the Gini coefficient and the profit share as a percentage of gross domestic income. What makes the series analytically distinctive is not only its length but its theoretical openness: Rodriguez-Weber explicitly recognises coercive distribution --- the state shaping market income \textit{``by sending the police or military forces to dissolve a strike, or by pressuring the entrepreneurs to accept a rise in wages''} --- moving distributional analysis beyond market mechanisms to state power and class struggle as determinative forces. The long-wave finding is equally significant: Chilean inequality follows Kuznets-like waves driven by globalisation, institutional change, and class power, and the UP sits at the trough of a long inequality-reduction cycle beginning in the 1930s --- contextualising the Allende government not as aberrant but as the culmination of a decades-long distributional shift. \citet{JaimovichToledo2021} trace the causal chain connecting agrarian reform intensity during the 1962--1973 period to the persistence of Mapuche territorial conflict in the post-dictatorship era (1990--2016), using geocoded plot-level data and an instrumental variable strategy to identify a mechanism the dependency tradition named structurally but could not trace empirically. The Badia-Miró, Llorca-Jaña, and Salazar-Pinto cluster contributes empirically driven data reconstruction on prices, trade, and living standards that enriches the quantitative foundation available for long-run analysis. These contributions demonstrate what the field's methods can achieve when applied with care to the Chilean case --- and they provide part of the empirical foundation that this essay's own framework requires.

But Rodriguez-Weber's framework remains institutional-descriptive despite its data sophistication. Functional income distribution across classes is measured, but not modelled. The UP receives one paragraph jumping from political choice to polarization without engaging the economic terrain --- no profit squeeze, no balance-of-payments crisis, no wage-price dynamics. There is no aggregate demand function, no investment equation, no external constraint. The limitation is not idiosyncratic --- it reflects a structural feature of the field. HPE's best empirical work inherits its macro theory from elsewhere, and what it inherits is the neo-structuralist framework whose limitations the previous move exposed. \citet{AldunateGonzalezPrem2022}'s study of U.S. banking covert action during the Allende government is the most revealing exhibit: a methodologically rigorous difference-in-differences design isolating the causal effect of U.S. credit withdrawal on Chilean firms, yet its entire macro framing is compressed into a single paragraph deferring to \citet{LarranMeller1991} --- the same populist macro paradigm whose strategic omissions the previous move documented. A paper about an external credit squeeze during a period of explosive distributional conflict takes the macroeconomic dynamics as given rather than modelling how the credit withdrawal interacted with profitability, the wage share, and the balance of payments as endogenous variables. The pattern is not idiosyncratic. \citet{GonzalezPrem2026} defer to Dornbusch and Edwards with the distancing word ``believed'"; \citet{GirardiBowles2018}, despite Bowles's heterodox credentials, cite the same NST references in a footnote. Across the entire HPE corpus reviewed for this chapter, zero of sixteen sources cite Kalecki, Marx, Shaikh, Thirlwall, Prebisch, or the surplus approach. The exclusion is therefore not a gap in coverage --- it is a boundary condition of the field's disciplinary constitution.

This is not a criticism of individual authors' competence --- it is a diagnosis of what the field makes available. \citet{Gailmard2021} identifies the problem from within: HPE research is empirically driven and undertheorised, prioritising the identification of causal effects over the development of theoretical frameworks capable of organising those effects into a coherent account of historical dynamics. \citet{Rozenas2021} concurs, noting that the field leverages mostly exogenous variation --- consequences of nature, historical accidents, policy discontinuities --- leaving history as a vacuum within which the identified causal relationship takes place. The credibility revolution's methods are designed to identify \textit{effects}; they are not designed to specify \textit{mechanisms} that connect those effects to the structural dynamics that produced the historical situation under study.

The deeper issue is disciplinary. \citet{JenkinsRubin2024} present HPE as an interdisciplinary intersection --- a Venn diagram where history, political science, and economics overlap. The self-description is accurate as \textit{soft-institutional sociology}. But the Venn diagram reveals what it does not intend to. The economics that enters the intersection is neoclassical economics --- and, where macro is required, whatever framework the discipline's current consensus provides. Classical political economy --- the tradition that generated the categories of surplus, accumulation, class, and crisis as analytical tools for understanding historical dynamics --- constitutes a different intersection entirely. HPE does not occupy it. It is not that the field has considered classical political economy and found it wanting. It is that the field's disciplinary constitution excludes it before the question arises.

HPE does not lack theory through oversight. It has a theory --- the neoclassical synthesis it inherits from its parent disciplines, transmitted through graduate training, reproduced through journal editorial practice, and reinforced by the methodological commitments that define what counts as admissible evidence. The exclusion of classical political economy is not a gap waiting to be filled. It is the field's constitutive boundary. What matters is not whether HPE's methods are capable --- Rodriguez-Weber's distributional series, the Aldunate identification strategy, Jaimovich and Toledo-Concha's geocoded agrarian reform instrument demonstrate that they are --- but what theoretical container those methods operate within. The questions the dependency tradition raised --- how the external credit squeeze interacted with profitability, how distributional conflict endogenised macroeconomic dynamics, how the terms of trade set a structural corridor within which class forces contend --- are questions that HPE's methods could address with considerable rigour under a different framework than the one the field currently permits. What this essay offers is therefore not an invitation to HPE to supplement its toolkit. It is a demonstration that the dependency research programme, renewed with mechanistic theorising \citep[pp.~15--19]{MadariagaPalestini2021}, provides the framework those methods require --- and that framework operates outside the boundary HPE's disciplinary constitution cannot cross. The gap between what HPE's methods can do and what its inherited theory allows it to ask is not a deficit that engagement will resolve. It is a structural condition that this essay addresses by building outside the field's boundary. If no existing school formally connects external transmission mechanisms, internal distributive conflict, and macroeconomic performance as endogenous, interacting dynamics --- the circle returns to where it opened, but now equipped with the tools the classical tradition lacked.

% ============================================================
% S1.4 — MODIFIED: s_pi = 1.381 insertion
% COSMIC LETTER: Discovery 3 — zeta_3 two dimensions; s_pi finding
% ============================================================

\subsection{Historical Political Economy --- the UP distributive conflict as empirical datum}
\label{sec:A4b_HP_UP_datum}

% COSMIC LETTER: Discovery 3 — s_pi = 1.381 from Ch2 data analyst
Before the circle closes with the revival, one empirical finding from
Chapter~2 must be placed. The structural break coefficient
$s_{\pi} = 1.381$, estimated in the Chile Stage~B Weisskopf
decomposition, establishes that the 1972--74 episode was distributionally
unique in both the US and Chilean series---the only period in which the
profit-share break exceeds the output break. What makes the finding
structurally significant is not its magnitude alone but its
two-dimensional character: the distributive conflict of the Unidad
Popular period was self-undermining through \textit{both} the
profit-squeeze channel \textit{and} the financial fragility channel
simultaneously. Rising $\omega$ compressed profit margins directly
(the standard Goodwin mechanism), while in an economy without
wage-goods sovereignty, the same rise in $\omega$ increased import
pressure through the social reproduction circuit---wage-goods demand
that could only be met through imports drained the reserves that
sustained the leverage ratio. The net effect, captured by $\hat{\zeta}_3
= -3.919$ as a semi-elasticity (approximately 3.9\% import reduction per
percentage-point wage-share increase; approximately 32\% per 10
percentage points), conceals an internal contradiction between
luxury-goods import compression (a source of relief) and wage-goods
import pressure (a source of fragility). This two-dimensional character
is what made the 1972--74 episode structurally unique: distributive
conflict at the periphery carried a financial fragility dimension that
the centre did not share, because the centre's wage goods were
domestically produced. The significance for what follows is direct: when
the leverage ratio is formalised in Section~\ref{sec:C_levr}, the wage
share enters through both channels, and suppressing either
misrepresents the peripheral dynamics.


\subsection{The Revival of Dependency Theory as a \textit{Research Programme}}
\label{sec:A5_rev_deptheory}

% [UNCHANGED — see original, lines 171–193]

The return to dependency theory is not a retreat. The tradition documented in the opening move of this chapter suffered a double defeat --- political and intellectual --- but \citet{Palma2016} insists on a distinction the subsequent literature collapsed: the dependency programme was not refuted, it was silenced. Neoliberalism's conquest of Latin America's progressive intelligentsia operated not through superior argument but through the elimination of institutional infrastructure --- the CESO dismantled, ESCOLATINA closed, the ECLAC tradition sanitised by the technopols who would go on to administer the post-dictatorship order. TINA functioned as a self-fulfilling prophecy: by destroying the spaces that could generate alternatives, the neoliberal victory became self-reinforcing. The arguments that had explained the UP's structural conditions --- the centre-periphery framework, the external constraint, the endogeneity of distributive conflict --- were never answered. They were made organisationally impossible to pursue.

The epistemological consequence is the reframing \citet{Kvangraven2021} provides. Applying Lakatos's methodology of scientific research programmes, she reconstructs dependency theory as \textit{progressive} --- a tradition that continues to generate novel predictions and empirical findings --- against the standard dismissal that it had become degenerative or merely ideological. The programme's hard core holds: global capitalism tends to be polarising. Its methodological commitments --- reading economic dynamics through a global historical lens, tracking how capitalism produces and reproduces peripheral structures of production, attending to the external constraints that set the corridor within which domestic accumulation can proceed --- describe with precision what this essay's empirical core undertakes. These are not abstract epistemological positions. They are research directives. The dependency tradition was not superseded by superior theory. It was interrupted before the work was done.

The programme's Lakatosian credentials are reinforced from a second disciplinary direction. \citet{MadariagaPalestini2021} formalise the dependency research programme around three core assumptions: (1) the global political economy is hierarchically structured into core and periphery; (2) peripheral dynamics are conditioned by core dynamics; and (3) the internal structure of peripheral units reflects the type of relation between periphery and core \citep[pp.~7--10]{MadariagaPalestini2021}. The third assumption is the one Part C's empirical strategy operationalises directly: Chile's distributional dynamics --- wage share, profit rate, capacity utilisation --- as shaped by the external corridor set through commodity terms of trade and the balance-of-payments position. Where Kvangraven provides the epistemological scaffolding, Madariaga and Palestini provide the analytical architecture: the distinction between \textit{situations of dependency} --- the concrete, historically bounded manifestations of dependent capitalism --- and \textit{mechanisms of dependency} --- the causal processes through which the conditioning of peripheral units by core dynamics actually operates \citep[pp.~15--19]{MadariagaPalestini2021}. This distinction directly addresses the diagnostic that the previous move applied to HPE. The credibility revolution developed powerful tools for causal identification but no framework specifying the \textit{mechanisms} that connect external constraint to internal dynamics. The dependency research programme now provides both the vocabulary and the empirical template --- mechanisms are not merely invoked but formally specified and traceable through the data. That the credibility revolution's claim to atheoretical causal identification is itself contested --- the selection of exogenous variation is theory-laden, and design-based approaches systematically privilege identification over mechanism specification --- only strengthens the case for a theoretically grounded alternative.\footnote{The critique of the credibility revolution's methodological premises, including its claims to design-based inference: [REFERENCE MISSING PLACEHOLDER]} Cardoso himself had anticipated the methodological point: criticizing Chase-Dunn's attempt to test dependency through large-N cross-country regression, he insisted that he was not arguing against the use of empirical data for validation or rejection of theories but against their inadequate use --- against a methodology that imposes formal and a-historical comparisons where the specificity of concrete situations is a precondition for analytical formulation \citep[p.~11]{MadariagaPalestini2021}.

The work has resumed --- and it resumed, first, as archaeology. \citet{CardenasLanaSeabra2022} is the single most concentrated evidence of what Palma diagnosed: a tradition not refuted but buried. Their edited volume \textit{unearths} dependency theory --- reproducing original research blueprints that had circulated in Santiago in the months before the coup and then disappeared into exile, and reassembling insider testimonies from \textit{dependentistas} members after decades of enforced silence. The tradition was not lost because it failed. It was physically removed. What is recovered is not merely an analytical programme but a normative one: a commitment to development as the overcoming of economic subordination, not its management. \citet{Ahumada2023} names the principle: freedom as non-domination, applied to centre-periphery relations, a normative commitment that lived within the tradition, but reintroduced under the light of contemporary political-philosophy of peripheral development. The unearthing restores it. What it cannot, alone, provide is the macro grounding that gives it analytical purchase --- and that is precisely what the revival's tools now supply.

But the revival is not only archaeological. The dependency research programme is being applied to contemporary capitalism in ways the classical tradition could not have anticipated. \citet{Ruiz2021} analyses Chile's privatised pension system as a mechanism of dependent financialisation: the Administradoras de Fondos de Pensiones function as a ``financial valve'' mediating local and transnational accumulation, channelling workers' social security contributions into the portfolios of national and international financial groups \citep[pp.~231--234]{MadariagaPalestini2021}. This is Cardoso's associated-dependent development applied to twenty-first-century Chile --- the internal class structure shaped by the external relation, as the programme's third core assumption specifies.

The comparative evidence is equally substantial. \citet{Bizberg2021} maps the diversity of dependent capitalisms across contemporary Latin America --- international outsourcing, socio-developmental, and two variants of rentier capitalism --- integrating regulation theory with the dependency framework and demonstrating that the tradition generates comparative typologies with genuine explanatory power \citep[pp.~131--155]{MadariagaPalestini2021}. \citet{BurchardtDietzWarnecke2021} theorise rent as a mechanism of dependency: primary commodity rents appropriated locally and realised on global markets, structurally conditioning the wage-profit relationship from the outside \citep[pp.~207--229]{MadariagaPalestini2021}. These are not historical reconstructions of what the tradition once said. They are new applications --- the programme generating novel findings in domains the classical version never reached, which is precisely what Lakatos' criterion of progressiveness requires.

What the classical tradition correctly identified, it could not formalise. The closing assessment at \S\ref{sec:A1_classical_deptheory} was precise: there was no theoretical unity across the tradition niether a formal specification of the transmission mechanisms connecting the external constraint to internal distributive dynamics. The revival does not resolve these limitations but addresses the one two interconnected tools now exist on the external side where the classical tradition had categories without models. The first is the formalisation of the external constraint itself. Thirlwall's balance-of-payments constrained growth model provides what \citet{Prebisch1981}'s mature structuralism had approached analytically but left without dynamic specification: a formal account of how the terms of trade determine the corridor within which domestic accumulation can proceed. Favourable terms open the corridor loosening the external constraint, enabling wage growth to compensate profit loses through demand channels. Once deteriorating terms close reverse, distributive conflict and rising consumption demand from abroad, explosive regardless of domestic policy choices.


The structural conditions documented in Chapter 2--- the peripheral economy's exposure to import of machinery as constraint for the expansion of productive capacities through mechanization --- ground this mechanism in the Chilean case specifically. This is what the classical tradition identified narratively. What Thirlwall's model specifies as a balance-of-payments corridor, \citet{Leiva2008} had diagnosed as a political economy: the terms of trade as the form through which peripheral surplus is appropriated and withheld, not merely the rate at which imports and exports clear. \citet{TorresAhumada2022} supply the mechanism --- three forms of rentierism structurally conditioning the wage-profit relationship from the outside --- that converts Leiva's diagnosis into a testable empirical claim. The LP chain maps onto all three layers simultaneously: the corridor (Thirlwall), the political economy of the corridor (Leiva), and the monetary mechanism through which the corridor operates (Section B) --- which is what makes Part C theoretically richer than a balance-of-payments specification alone. \citet{TorresAhumada2022} renew the centre-periphery framework not as historical description but as a theory of how primary commodity rents are appropriated, circulated, and withheld under contemporary peripheral capitalism --- three forms of rentierism (classic resource extraction, financial, digital) that restore the surplus approach \citet{Leiva2008} diagnosed NST as having jettisoned, and provide the macro foundation that the HPE saturation check confirmed the field entirely lacks. The LP framework tracing U.S.\ military-Keynesian expenditure through copper terms of trade to domestic growth-distribution dynamics is precisely one instance of centre-periphery rentierism as \citet{TorresAhumada2022} theorise it --- the essay's own contribution is legible within this framework without requiring a new theoretical apparatus.

The external tools specify the corridor; the internal tool specifies what happens within it. \citet{Sossdorf2024} documents Latin American productivity-wage dynamics from 1975 to 2018, establishing that Chile's positive decoupling coefficient --- productivity gains do transmit to wage growth, at a ratio of 1.73 --- is not the region's norm but requires explanation. Within a Kaleckian framework, it becomes explicable: the wage share and profit rate are endogenous outcomes of the interaction between the external conditions set by the BoP position and the domestic class configuration. Distribution is not a policy parameter set outside the model --- it is produced inside it, by the encounter between structural corridor and class power. This is precisely what the classical tradition had the ambition to theorise but never the instruments to formalise, what NST ruled out by construction when it exogenised distribution, and what HPE cannot recover because it has no macro theory of its own. The original tradition had the ambition. The revival has the instruments.

Those instruments have a precise empirical deployment. \citet{Kvangraven2021}'s four methodological tenets --- global historical approach, theorising polarising tendencies of capitalism, focus on structures of production, attention to peripheral constraints --- describe with near-exact correspondence what Part C's local projections event study undertakes for Chile across 1925 to 2010. The first link in the LP chain traces U.S. military-Keynesian expenditure as an exogenous demand shock propagating through copper terms of trade to Chile's external position: centre-country fiscal-military policy setting the peripheral structural corridor, as the renewed centre-periphery framework specifies. In the vocabulary Madariaga and Palestini provide, this LP chain is itself a \textit{mechanism of dependency} --- the concrete causal process through which the dynamics essential for the maintenance of Chile's political-economic order are conditioned by the dynamics of the U.S. fiscal-military apparatus. The monetary mechanism through which this propagation operates --- the hierarchical structure of world money and the peripheral state's subordinate position within it --- is theorised in the essay's empirical section, drawing on Marx's category of world money and its contemporary elaboration by \citet{Basu2022}. The second link traces that corridor's domestic consequences --- capacity utilisation, wage share, profit rate, international reserves, private investment --- each variable endogenous, each channel formally specified, conditioned on the BoP position that determines whether the corridor is open or closing. The post-1971 narrowing that preceded the coup becomes legible in this framework not as policy error, not as socialist mismanagement, and not as the inevitable consequence of a reformist strategy --- but as the interaction of an externally induced constraint with an internally explosive distributive dynamic that no school reviewed in this chapter had the theoretical equipment to model simultaneously. The chapter opened with a tradition whose theoretical foundations were still being assembled when the coup interrupted them. It ends with the same tradition --- not restored but extended, not repeated but equipped. What the Unidad Popular's political-economic programme demanded from theory, and did not receive in time, can now be formally given: a connected account of how the external constraint becomes internal distributive crisis, and how that crisis, in turn, becomes history.

% ============================================================
% S1 CLOSE — TRANSITIONAL PARAGRAPH (NEW)
% COSMIC LETTER: Discovery 1 — dependency renewed needs world-money formalization
% ============================================================

% COSMIC LETTER: Discovery 1 — peripheral-first reframing
The revival restores the dependency research programme as a progressive
Lakatosian enterprise and equips it with the formal tools---Thirlwall's
BoP-constrained growth model, the Goodwin--Kalecki distributive cycle,
the Shaikh profit-rate decomposition---that the classical tradition
lacked. But the programme's completion requires one further step that
no contribution reviewed in this section provides: the formalisation of
world money as the monetary architecture through which the external
constraint operates. The balance-of-payments corridor is not merely a
flow constraint on trade; it is a balance-sheet position of the
peripheral central bank denominated in a currency the peripheral economy
does not issue. What converts the technological dependency documented in
Chapter~2---the mechanisation trap, the structural import
propensity---into the financial fragility the UP confronted is the
monetary hierarchy through which peripheral accumulation is subordinated
to the centre's currency. That monetary architecture, and the leverage
ratio it produces, is the object of the next section.


% ============================================================
% SECTION 2 — FINANCIAL SUBORDINATION AND WORLD MONEY
% Modifications: S2.3 (Bretton Woods asymmetric transmission opening)
%                S2.4 (zeta_3 financial fragility connection)
% ============================================================

\section{Financial Subordination and World Money}
\label{sec:B_FS_WM}

\subsection{Historical transit: commodity money to credit money}
\label{sec:B1_historymoney}

% [UNCHANGED — see original, lines 198–205]

Money does not originate as credit. It originates as the commodity form taken by the universal equivalent---the one commodity whose use-value is to be the social expression of the value of all others \citep[\S3.1]{Basu2022}. In Marx's account, the exchange process itself produces a single commodity whose social function is to serve as the general measure of value and the universally accepted means of payment. Gold occupies this position not because of a natural property but because a historically specific process of social validation singles it out: the function is socially determined, the form is historically contingent \citep[\S3.2.1]{Basu2022}. What matters for theoretical purposes is not gold but the universal equivalent as a category---the socially validated device through which privately produced commodities are recognised as bearers of abstract labour.

Marx traces the developmental sequence by which credit money emerges from this commodity foundation, but the analysis remains at a level of abstraction he did not fully close. \citet[\S3.2.3]{Basu2022} reconstructs the chain. The bill of exchange, issued in the course of commercial transactions, is discounted by a bank; the bank issues its own liability against the discounted bill; the central bank stands behind commercial bank liabilities as the bank of banks; the state makes central bank liabilities legal tender. At each step the credit instrument substitutes for the commodity anchor, and at the final step social validity---the condition any form of money must satisfy to function as the universal equivalent---is guaranteed by the institutional authority of the state, not by the intrinsic value of a commodity \citep[\S3.2.3]{Basu2022}. The commodity anchor is no longer needed. What has changed is the form; what has not changed is the social function. Central bank liabilities and commercial bank deposits together now perform what gold once performed: the social materialisation of value, the mechanism through which privately produced commodities are validated as socially necessary.

The operative conversion factor linking value magnitudes to monetary magnitudes in a credit money system is the monetary expression of value---the ratio at which abstract labour-time is expressed in the monetary unit \citep{Foley1982, Basu2022}. What remains to be established is what this means when the universal equivalent crosses the border.


\subsection{World money under the Monetary Expression of Value }
\label{sec:B2_worldmoney_mev}

% [UNCHANGED — see original, lines 209–219]

Marx's world money is bullion. When money leaves the sphere of national circulation it strips off its local garbs and returns to its original form---the commodity that functions as the absolute social materialisation of wealth \citep[p.~51]{Ivanova2013}. In Marx's formulation this is not a secondary function but the fullest expression of the potentialities and contradictions of the money-form: in the world market, money functions as the commodity whose natural form is the direct social form of realisation of human labour in the abstract. The question \citet{Ivanova2013} poses is whether this category can extend to a national fiat currency---whether a credit money instrument issued by a particular state can perform the function that Marx reserved for bullion.

Her answer turns on a distinction between substance and function. The fundamental feature of Marx's theory, she argues, is not its commodity-basis but its conception of money as the social expression of value, whose magnitude is determined by socially necessary labour embodied in the commodity. Money is an expression of the class relation between labour and capital, and the management of money is inseparable from the social form of imposition of work \citep[p.~45]{Ivanova2013}. If this is the theoretically load-bearing condition---not the physical substance of the universal equivalent but its social function---then the form of the universal equivalent is contingent. On this reading, the dollar is the closest equivalent to genuine world money the capitalist world market has known \citep[p.~52]{Ivanova2013}, not because it is a commodity but because it performs the function: it is the universally accepted means of international payment, reserve asset, and unit of account.

\citet[p.~6]{Ivanova2020} specifies the condition under which this claim holds. Credit money becomes money---income rather than mere promise---only if it enters the domain of social production as the embodiment of the value of social labour and social purchasing power. Central bank liabilities acquire the attributes of the universal equivalent not by state decree alone but through their social validation in production and circulation. The dollar's world money status is therefore not a legal fact but a social relation, continuously reproduced through the global circuit of capital. This condition is what separates Ivanova's reading from a merely nominalist account of fiat money.

The monetary expression of value framework provides a lens through which to assess this argument. The monetary expression of value, defined as the ratio of monetary value added to productive labour hours, is the operative conversion factor within any national monetary system \citep{Foley1982}. The Marxian exchange rate---the ratio of one country's monetary expression of value to another's---preserves social labour-time in exchange: it is the rate at which one hour of productive labour in economy~A exchanges for one hour in economy~B \citep[\S3.2.7]{Basu2022}. Ivanova's form-function argument is coherent with this framework. The dollar as world money means the US monetary expression of value functions as the implicit international reference against which peripheral monetary expressions are measured. This is not a formal derivation from the framework---it is a compatibility that allows taking her account at face value for its central claim.

What Ivanova's account does not fully theorise is the mechanism by which this position is structurally reproduced. Her analysis establishes that the dollar performs the world money function and that this performance is grounded in social relations of production, not in commodity substance or state fiat. But it does not explain why the dollar specifically retains this position, nor why peripheral subordination deepens rather than dissolves as US industrial dominance erodes. Ivanova identifies the explanandum; the explanans requires the institutional and structural account that follows.


% ============================================================
% S2.3 — MODIFIED: Bretton Woods opened with Ch2 asymmetric transmission
% COSMIC LETTER: Discovery 4 — 1972–74 asymmetric transmission as empirical datum
% ============================================================

\subsection{Bretton Woods and the peripheral balance-of-payments constraint}
\label{sec:B3_BrettonWoods_BOP}

% COSMIC LETTER: Discovery 4 — Bretton Woods asymmetric transmission as empirical datum from Ch2
Chapter~2's four-channel Weisskopf decomposition established an empirical
asymmetry that demands a monetary explanation. The 1972--74 contraction
was price-mediated at the centre: 57.8\% of the US profit-rate decline
operated through the $P_Y/P_K$ channel---a one-sided shock absorbed
through relative prices within a monetary system the United States itself
governed. For Chile, the same conjuncture produced a two-sided shock:
the collapse of the copper price compressed export revenue while the
dollar float---the dissolution of the gold-dollar peg---simultaneously
destabilised the exchange-rate anchor of Chile's reserve position. The
question this asymmetry poses is structural: what monetary architecture
produces a shock that is one-sided at the centre and two-sided at the
periphery? The answer is the gold-dollar standard itself, and the
leverage ratio formalises the transmission mechanism through which the
asymmetry operates.

The dollar-gold peg at \$35 per ounce, operative from 1944, institutionalised the dollar as the form of world money during the essay's period of interest. The United States occupied the apex of the monetary hierarchy not through market competition but by institutional design: it emerged from the war as the dominant creditor with substantial gold reserves, and the Bretton Woods architecture codified this position into the operating rules of the international monetary system \citep[p.~483]{Vasudevan2009}. The exorbitant privilege---the capacity to borrow internationally in one's own currency and to earn returns on foreign assets that exceed the cost of foreign liabilities---was operative from the outset, not only after the system's collapse \citep[p.~473]{Vasudevan2009}.

The asymmetry between centre and periphery under this architecture is structural, not contingent on any particular balance of trade. The country issuing world money can finance current account deficits in its own currency, and in doing so it centralises and recycles international surpluses through its financial system \citep[pp.~243--244]{Vasudevan2019}. It acts, in effect, as the banker to the world---borrowing short in the form of reserve holdings by other central banks and lending long in the form of foreign investment and military expenditure. Peripheral countries cannot do this. Their external liabilities are denominated in the dominant currency, and the settlement of international payments requires access to that currency. Instead of being able to shift the burden of adjustment onto trading partners, as the dominant country can, the periphery bears the brunt of debt-deflationary dynamics disproportionately \citep[p.~486]{Vasudevan2009}.

The periphery therefore operates under a balance-of-payments-constrained circuit of capital \citep{Vasudevan2016, Abeles2013}. Domestic credit creation is semi-endogenous---not in the sense that the central bank has no discretion over monetary policy, but in the sense that the ceiling on sustainable domestic monetary expansion is set by the accumulation of reserves denominated in the dominant currency, by export earnings, and by the terms on which external borrowing is available. The peripheral central bank holds the dominant country's liabilities as the necessary form of international liquidity, with no symmetric obligation on the other side. Reserve accumulation under these conditions is not a policy choice. It is the structural form through which the peripheral monetary system is anchored to the international monetary hierarchy---a form of financial subordination that reproduces the asymmetry regardless of domestic policy orientation \citep[p.~244]{Vasudevan2019}.

Chile's monetary system during the Bretton Woods period operates within this constraint. The closing of the gold window in 1971 dissolved the nominal anchor without dissolving the dollar's world money function \citep[p.~486]{Vasudevan2009}---a structural break whose implications for the transmission mechanism are tested through the post-1973 placebo specification in Section~\ref{sec:empirical_strategy}.


% ============================================================
% S2.4 — MODIFIED: zeta_3 financial fragility dimension connected
% COSMIC LETTER: Discovery 3 — financial subordination as structural condition
%   that makes the wage-goods import dependency channel operative
% ============================================================

\subsection{Financial subordination as a state variable}
\label{sec:B4_fs_levR}

The notion of world money, examined in the preceding subsections through
the commodity-to-credit transition, the form-function argument for the
dollar, and the Bretton Woods architecture of peripheral subordination,
finds its operative expression in the peripheral central bank's balance
sheet---the ratio of its obligations to its own population to its claims
on monetary value from the hegemon of the capitalist world system. What
remains is to specify the forces that govern its movement.

% COSMIC LETTER: Discovery 3 — zeta_3 two-dimensional character maps onto
%   financial subordination as the structural condition for wage-goods import dependency
The connection to Chapter~2's empirical findings is direct. The
coefficient $\hat{\zeta}_3 = -3.919$, estimated as a semi-elasticity in
the Chile TVECM, carries a two-dimensional character that the single
parameter conceals. Through the standard profit-squeeze channel, rising
$\omega$ compresses $\pi$, slows accumulation, and reduces import demand
for capital goods---a relief channel. Through the financial fragility
channel, rising $\omega$ in an economy without wage-goods sovereignty
raises import pressure through the social reproduction circuit: the
consumption basket of the working class contains imported wage goods
whose demand rises with the wage share, draining the reserves that
sustain the leverage ratio. Financial subordination---the structural
position formalised in this section---is precisely what makes this second
channel operative. An economy with wage-goods sovereignty (the US case)
absorbs distributional gains domestically; an economy subordinated within
the monetary hierarchy (Chile) transmits those gains into reserve
drainage through the import channel. The leverage ratio, constituted in
the next section, is the balance-sheet object through which this
transmission is registered and through which it becomes a binding
constraint on accumulation.


% ============================================================
% SECTION 3 — FORMALIZATION OF THE LEVERAGE RATIO
% Modifications: S3.2 (empirical anchors from Ch2), S3.3 (bridge to S4)
% ============================================================

\section{Formalization of the Leverage Ratio}
\label{sec:C_levr}

\subsection{The leverage ratio as a balance sheet object}
\label{sec:C1_levr_constitution}

% [UNCHANGED — see original, lines 269–357]

The monetary architecture established in the preceding section identifies the
peripheral central bank's balance sheet as the institutional site where
financial subordination takes its operative form. The leverage ratio of the
peripheral central bank is not a reserves adequacy metric. The mainstream
open economy literature frames international reserves as a policy instrument
to be optimised---a buffer the central bank sizes against financial
frictions, choosing the stock that maximizes welfare under the constraint
that reserves are costly to hold.\footnote{See, e.g.\ a seminal article by
	\citet{Jeanne2007} and a more recent contribution for open small economies
	like Chile by \citet{CespedesChang2024}.}
The question that literature asks---how large should the buffer
be?---presupposes what this essay
contests: that reserves accumulation is a policy choice rather than a
structural position. In the framework developed here, the leverage ratio is
not a quantity the central bank optimises. It is a balance sheet object
whose two sides are denominated in theoretically distinct forms of money,
and what it encodes is not a prudential decision but the structural
coordinate of the peripheral monetary system within the international
monetary hierarchy.

On the liability side, currency in circulation---$M_t$---is an obligation
of the central bank to the population: the final link in the credit money
chain that \citet[\S3.2.3]{Basu2022} reconstructs, in which the state's
institutional authority guarantees social validity without a commodity
anchor. The peso in a Chilean worker's pocket is the central bank's
debt---the domestic universal equivalent issued as a passive of the
institution that stands at the apex of the national monetary system. On the
asset side, international reserves---$IR_t$, denominated in US
dollars---are a claim on monetary value expressed in the currency that
performs the world money function. At the level of abstraction established
in \S\ref{sec:B2_worldmoney_mev}, which currency occupies this position is
contingent; that some currency must occupy it is a structural requirement of
the world market under credit money. But the dollar is not an abstract
placeholder. It is the concrete expression of financial subordination to
the hegemon of international transactions---a country that imposed its own
currency as world money through the architecture of the hierarchical order
of the inter-state system \S\ref{sec:B3_BrettonWoods_BOP} documented, and
that reproduces its position of power, not through value equivalence.
The ratio of the two sides is therefore a leverage ratio in the balance
sheet sense:
%
\begin{equation}
	\label{eq:levr_def}
	\widetilde{\text{LevR}}_t \;\equiv\; \frac{e_t \cdot IR_t}{M_t}
\end{equation}
%
where $e_t$ is the nominal exchange rate (domestic currency per dollar).
What this ratio measures is the central bank's capacity to conduct monetary
policy under conditions of financial subordination---a capacity that is
real, but whose terms are set by the hierarchical position of the national
monetary system within the international monetary order. The normative
question---how large should the buffer be?---dissolves once the object is
reconstituted: the leverage ratio is not a buffer to be sized but a
structural coordinate to be read.

The growth-rate decomposition names the mechanism:
%
\begin{equation}
	\label{eq:levr_growth}
	g_{\widetilde{\text{LevR}},t} \;=\; g_{e,t} \;+\; g_{IR,t} \;-\; g_{M,t}
\end{equation}
%
Three structurally distinct forces act on the leverage position. Exchange
rate depreciation revalues the asset side in domestic currency terms---a
pure balance sheet effect. Balance-of-payments flows add or drain
reserves---the external channel through which terms-of-trade shocks, debt
service, and capital movements reach the central bank's asset position.
And domestic credit creation expands liabilities---the internal channel
through which fiscal monetisation, wage-price dynamics, and the distributive
conflict itself erode the central bank's leverage from below. What matters
is not only whether the asset side is deteriorating but whether liability
expansion is outpacing it: the central bank can erode its own capacity to
conduct monetary policy by expanding domestic currency issuance even when
reserves are stable. Under the gold standard, this ratio had a direct
physical expression---gold reserves against notes in circulation, with the
leverage position materially constrained by the statutory reserve
requirement. The credit money reconstruction in
\S\ref{sec:B1_historymoney} makes legible what changed and what did not:
the commodity anchor dissolved domestically, but the Bretton Woods
architecture preserved the analogous balance sheet structure
internationally. The gold coverage ratio was always a leverage ratio.
$\widetilde{\text{LevR}}_t$ names the same structural object in its
credit-money form. The central bank does not choose its leverage---it
inherits it from the interaction of external flows that govern the
numerator and domestic monetary expansion that governs the denominator.
The ratio has been constituted; what governs its dynamics is the
balance-of-payments identity.


% ============================================================
% S3.2 — MODIFIED: empirical anchors from Ch2 inserted
% COSMIC LETTER: zeta_1 = 0.930, gamma_hat = -0.1394
% ============================================================

\subsection{Balance-of-payments dynamics of the leverage ratio}
\label{sec:C2_bop_dynamics}

The leverage ratio has been constituted as a balance sheet object; its dynamics
have not. What governs how $\widetilde{\text{LevR}}_t$ moves is the
balance-of-payments identity---the accounting relation that determines the
flow into and out of the central bank's asset side. The general
balance-of-payments equation states that the change in the peso value of
international reserves equals the sum of the current and capital account
balances:
%
\begin{equation}
	\label{eq:bop_identity}
	\Delta(e_t \cdot IR_t) = CA_t + KA_t
\end{equation}
%
where $CA_t$ carries the trade balance, debt service, and net transfers,
and $KA_t$ carries net external borrowing, foreign direct investment, and
capital flight.\footnote{The full decomposition of $CA_t$ and $KA_t$ into
	their component terms, the substitution algebra, and the constraint-binding
	cases under alternative scenarios are developed in
	Appendix~\ref{app:bop_levr}.} Three structurally
distinct channels operate through this identity: the trade channel, through
which terms-of-trade movements compress or expand export revenue relative
to imports; the financial channel, through which debt service and capital
flight drain reserves independently of trade; and the valuation channel,
through which exchange rate depreciation revalues the existing stock of
dollar reserves in domestic currency---a pure balance sheet effect that
changes $\widetilde{\text{LevR}}_t$ without altering the underlying real
position.

Substituting the leverage ratio definition $e_t \cdot IR_t =
\widetilde{\text{LevR}}_t \cdot M_t$ into
equation~\eqref{eq:bop_identity} and rearranging yields the law of motion
for the leverage ratio:
%
\begin{equation}
	\label{eq:levr_lom}
	\Delta\widetilde{\text{LevR}}_t \;=\;
	\frac{CA_t + KA_t}{M_t}
	\;-\; \widetilde{\text{LevR}}_t \cdot g_{M,t}
\end{equation}
%
The first term is the net balance-of-payments flow scaled by the domestic
monetary base---the rate at which external flows add to or drain the
central bank's asset side, expressed per unit of domestic liability. The
second term is the \textit{monetary expansion multiplier}: the existing
leverage ratio, scaled by the growth rate of domestic monetary liabilities.
This second term is what makes the leverage ratio a genuinely dynamic
object rather than a static coverage indicator. Even when the
balance-of-payments position is in surplus---when the first term is
positive---the leverage ratio falls whenever that surplus is insufficient
to cover the monetary expansion multiplier. The central bank's capacity
to conduct monetary policy erodes not only when reserves drain but when
liability expansion outpaces whatever reserve accumulation the external
position delivers.

The structure of the multiplier warrants attention. It is the
\textit{product} of two terms: the level of $\widetilde{\text{LevR}}_t$
and the growth rate of $M_t$. A central bank with a high leverage
position can sustain moderate monetary expansion without triggering the
constraint, because the multiplier threshold is high---the economy is far
from the binding region. A central bank with a low leverage position
faces a small multiplier in absolute terms, but the margin of safety is
also small---any monetary expansion, however modest, can push the ratio
below the critical threshold. The multiplier therefore operates
asymmetrically across leverage states: the same rate of monetary expansion
has different consequences depending on where the leverage ratio sits. This
asymmetry is what the state-dependent specification presented in
Section~\ref{sec:empirical_strategy} exploits as identification strategy
for the event study.

The leverage position is deteriorating whenever:
%
\begin{equation}
	\label{eq:binding_condition}
	\Delta\widetilde{\text{LevR}}_t < 0 \;\iff\;
	\frac{CA_t + KA_t}{M_t} < \widetilde{\text{LevR}}_t \cdot g_{M,t}
\end{equation}
%
A falling leverage ratio does not mean the constraint is binding---it means
the constraint is tightening, the distance to the binding threshold is
shrinking. Whether the constraint actually binds depends on the
\textit{level} of $\widetilde{\text{LevR}}_t$, not only on its direction
of movement. A central bank whose leverage is falling from a high position
retains policy space; a central bank whose leverage is falling from an
already low position may find that the same rate of deterioration
forecloses it. The constraint tightens in two structurally distinct ways. The
first is external: a terms-of-trade deterioration, a surge in debt
service, or capital flight compresses the numerator of the first term in
equation~\eqref{eq:levr_lom}---the balance-of-payments flow
itself turns negative, and reserves drain from the asset side. The second
is internal: the distributive dynamic operating through the domestic
economy raises the growth rate of monetary liabilities independently of the
external position. When the wage share rises---through the Kaleckian
bargaining channel that the empirical strategy formalises as the second
conditioning variable---the fiscal and monetary consequences follow:
wage-driven inflation, government expenditure to sustain employment and
redistribution, central bank credit to the banking system. Each of these
expands $M_t$, raising $g_{M,t}$, and the multiplier converts that
liability expansion into leverage deterioration. The external constraint
and the domestic distributive dynamic are therefore not independent
forces acting on the peripheral economy from different directions.
They meet in the leverage ratio, and the monetary expansion multiplier is
the mechanism of their interaction.

% COSMIC LETTER: zeta_1 = 0.930 and gamma_hat = -0.1394 as empirical special cases
The law of motion in equation~\eqref{eq:levr_lom} has empirical special
cases that Chapter~2's Chile estimation identified. The structural import
propensity $\hat{\zeta}_1 = 0.930$, estimated in the TVECM
cointegrating vector, establishes that each unit of machinery-and-equipment
accumulation adds 0.930 to the structural import requirement---a near
unit-elastic relationship that enters the first term of
equation~\eqref{eq:levr_lom} through the current account's import
component. The mechanisation trap documented in Chapter~2 is therefore not
merely a material constraint on the production structure; it is a direct
determinant of the rate at which the leverage ratio's asset side is
drained by the accumulation process itself. The TVECM threshold parameter
$\hat{\gamma} = -0.1394$ provides the empirical anchor for the
regime-switch condition: when the error-correction term exceeds this
threshold, the balance-of-payments dynamics enter the binding regime in
which import adjustment becomes endogenous to the constraint. This is the
empirical counterpart of the binding condition in
equation~\eqref{eq:binding_condition}---the point at which the leverage
ratio's deterioration forecloses the accumulation path that generated the
import requirement in the first place.

The balance-of-payments channel through which terms-of-trade shocks reach
the leverage ratio has an observable decomposition that the LP exploits as
a mechanism transparency device. The trade balance in nominal terms
responds to a terms-of-trade improvement on impact---the same export
volume commands a higher price, generating an immediate inflow of foreign
exchange. The trade balance in real terms---export and import volumes---responds
with a lag and potentially in the opposite direction: the income
effect of the terms-of-trade improvement triggers domestic demand
expansion, capacity utilisation rises, and import volumes surge through
the consumption and intermediate goods channels. The divergence between
the nominal and real trade balance impulse responses---what the empirical
strategy document names the \textit{trade scissors}---reveals the
mechanism through which the balance-of-payments constraint tightens from
the inside. The terms-of-trade improvement generates leverage headroom
through the nominal channel that is then consumed by the demand expansion
through the real channel. Whether the scissors is sustainable or explosive
depends on the leverage state: under high
$\widetilde{\text{LevR}}_t$, the central bank can absorb the import
surge without the constraint binding; under low
$\widetilde{\text{LevR}}_t$, the same import surge drains the remaining
margin and the constraint forecloses the demand expansion that generated
the surge in the first place. That is why the local projections
specification conditions the transmission of terms-of-trade shocks on
the leverage state---and why the formal apparatus developed here must be
complemented by the identification strategy that Section~\ref{sec:empirical_strategy} provides.


% ============================================================
% S3.3 — MODIFIED: bridge to Section D strengthened
% ============================================================

\subsection{From theory to identification}
\label{sec:C3_theory_to_id}

The preceding subsections established that the leverage ratio is not a
static indicator but a dynamic object governed by the balance-of-payments
identity, compounded by the monetary expansion multiplier, and subject to
a threshold condition that determines when the balance sheet forecloses the
central bank's policy space. The formal apparatus specifies a transmission
channel: terms-of-trade shocks affect the central bank's asset side through
the current account, and the monetary expansion multiplier translates the
domestic distributive dynamic into liability-side pressure on the leverage
position. Whether these two forces compound into a binding
constraint---whether the balance sheet forecloses policy space---depends on
the leverage state at the moment the shock arrives. The same terms-of-trade
deterioration that tightens the constraint under low leverage may be
absorbed without consequence under high leverage, because the multiplier
threshold is far from the binding region. What the apparatus does not
specify is what determines the \textit{distributional content} of whatever
passes through this channel when the constraint is not binding and the
corridor is open. The same terms-of-trade improvement, operating through
the same leverage position, can produce radically different domestic
outcomes depending on the balance of class forces at the moment the shock
arrives.

The unemployment rate is the variable that conditions this content. It
enters the empirical specification not as a labor market indicator in the
neoclassical sense but as a proxy for the balance of power between capital
and the working class---formal and informal, organised and unorganised.
The theoretical lineage is precise. The inverse relation between
unemployment and the wage share is Marx's reserve army mechanism: when
accumulation outpaces labour force growth, the reserve army thins, wages
rise, and the rate of exploitation falls. \citet{Goodwin1967} formalises
this as a dynamical system---the predator-prey cycle in which high
employment raises the wage share, compresses the profit rate, slows
accumulation, and restores unemployment, which disciplines wages back
down. What \citet{Kalecki1943} adds is the political dimension that
neither Marx nor Goodwin captures formally: low unemployment enables not
merely higher wages but labour \textit{militancy}---the organisational
and political capacity of the working class to press demands that exceed
the wage relation.\footnote{Kalecki's formulation posits a defined state
	of full employment. The argument here works without one: the mechanism
	intensifies continuously as the reserve army thins, which is why the
	specification partitions on a median split rather than a theoretically
	derived threshold.} It is this political intensification---strikes,
collective action, demands on the state for redistribution, the
organisational density of the working class as a class---that the
unemployment rate proxies for in the specification. When the reserve army
is thick and unemployment is high, the same demand expansion is absorbed
without distributional challenge: the surplus accrues to capital, and the
wage share remains disciplined. The leverage ratio governs the external
corridor; the unemployment rate governs what happens inside it. The local
projections specification in Section~\ref{sec:empirical_strategy}
therefore conditions on both state variables jointly: the leverage ratio
determines what is economically sustainable; the balance of class forces
determines what is politically demanded.

One scope condition must be stated explicitly. The structural level of
unemployment that obtained when the Unidad Popular took power in 1970---a
labor market already in the vicinity of full employment---was not
produced by the UP's own demand expansion. It was the product of a decade
of sustained loose external constraint through the long sixties: favourable
terms of trade kept the leverage ratio high, which permitted sustained
demand expansion, which absorbed the rural-urban labour flux through the
Okun channel and structurally thinned the reserve army. The unemployment
rate at time $t$ is therefore endogenous to the leverage ratio's trajectory
over the preceding decade. The local projections specification conditions
on both state variables contemporaneously---it cannot recover this
temporal hierarchy. The theory supplies the structural history that the
empirical design inherits as a given.

Three predictions follow from the formal apparatus and the class-power
argument jointly, each testable through the post-1973 placebo. First, the
wage share's response to a terms-of-trade shock should attenuate or
reverse after the coup, because labour repression dismantled the
bargaining infrastructure through which demand expansion translated into
distributional gains. Second, the leverage ratio's conditioning power
should weaken, because capital account liberalisation and foreign
borrowing altered the channels through which the balance-of-payments
constraint binds. Third, the unemployment rate's conditioning should lose
its distributional content---not because unemployment rose, but because
the political infrastructure through which low unemployment translated
into class power was destroyed. State terrorism enforced a reversal of
union density, but the lasting mechanism was institutional: the labour
code imposed under the dictatorship fragmented collective bargaining to
the firm level, dismantling the sectoral and inter-firm coordination
that had given organised labour its structural weight---a configuration
of union fragmentation that persists to the present day
\citep{Osorio2024}. Even when unemployment eventually fell again in the
1980s, the class-power content of low unemployment was structurally
different---workers had been disciplined by terror and then locked into
an institutional architecture designed to prevent the reconstitution of
class-wide bargaining power. If the transmission mechanism is structurally
specific to the pre-coup configuration, these three predictions should hold
simultaneously across the separately estimated chains.

% COSMIC LETTER: bridge to S4 strengthened
The formalization has produced two objects that the empirical strategy
operationalises directly. The first is the state variable itself: the
leverage ratio $\widetilde{\text{LevR}}_t$ as defined in
equation~\eqref{eq:levr_def}, whose law of motion
(equation~\ref{eq:levr_lom}) specifies how external flows and domestic
monetary expansion interact to determine the structural position of
the balance of payments. The second is the state-conditioning logic:
the $\{\text{LevR} \times u\}$ partition that combines the external
corridor (leverage ratio) with the internal balance of class forces
(unemployment rate) to produce a four-cell state space in which the
transmission of external shocks takes structurally distinct forms.
Section~\ref{sec:empirical_strategy} translates these two objects into
an estimable design: state-dependent local projections that condition
the impulse response of domestic outcomes on the leverage state, test
the regime-dependence null, and exploit the post-1973 structural break
as a placebo falsification of the mechanism's historical specificity.


% ============================================================
% SECTION 4 — EMPIRICAL STRATEGY (NEW — stacked from 07_section_D_draft_v1.tex)
% Adapted: preamble removed, equation numbering continues from S3,
%   cross-references to LevR definitions from S3, causal graph figure.
% ============================================================

\section{Empirical Strategy: State-Dependent Local Projections}
\label{sec:empirical_strategy}

This section formalises the empirical strategy used to estimate the causal chain
connecting Cold War military expenditure in the United States to domestic
macroeconomic outcomes in Chile during 1925--2010. The identification design
employs \textit{state-dependent local projections} \`{a} la \citet{Jorda2005},
organised around three empirical links. Link~1 estimates the reduced-form
external transmission from U.S.\ military spending to Chile's terms of trade.
Link~2 estimates the domestic propagation of terms-of-trade shocks into capacity
utilisation, distribution, and accumulation. Link~3 conditions Link~2 on the
balance-of-payments state variable to test whether transmission intensity is
regime-dependent. The causal architecture underlying these links is summarised in
Figure~\ref{fig:causal_graph}.

The design is anchored in Robinson's (1980) conception of \textit{Historical
Time}: causal compounding requires that upstream and downstream chains be
simultaneously active within the same calendar-year window. Statistical
association across non-overlapping periods carries no causal content. This
principle disciplines every stage of the estimation: the event-window
definitions, the state-dependent conditioning, and the compound multiplier
constructed in the results section all enforce temporal co-activation as a
necessary condition for causal inference.

% ==============================================================================
\subsection{Link 1: External Transmission}
\label{subsec:link1}

Link~1 identifies the reduced-form external transmission chain
$\Delta G^m \to \Delta \text{ToT}$, where $G^m_t$ denotes real U.S.\ federal
defence outlays and $\text{ToT}_t$ is Chile's net barter terms of trade. The
theoretical mechanism runs through Military Keynesian stimulus: defence
expenditure raises U.S.\ capacity utilisation ($\mu_{\text{US}}$), which
increases derived demand for industrial metals, bidding up the world copper
price ($P_{\text{Cu}}$), which in turn dominates Chilean export revenues. The
reduced form collapses this chain into a single estimable relationship.

The local projection specification for horizon $h$ is:
%
\begin{equation}
\label{eq:link1}
\Delta\text{ToT}_{t+h} = \alpha_h + \beta_h \,\Delta G^m_t
  + \gamma_h \,\widetilde{\text{LevR}}_t + \delta_h \,\Delta\text{USD\_Au}_t
  + \sum_{k=1}^{p} \rho_{h,k} \,\Delta\text{ToT}_{t-k}
  + \sum_{k=1}^{p} \varphi_{h,k} \,\Delta G^m_{t-k}
  + D^{\text{supply}}_t + \varepsilon_{t+h}
\end{equation}
%
where $\widetilde{\text{LevR}}_t$ is the Central Bank leverage ratio as defined
in equation~\eqref{eq:levr_def},
$\Delta\text{USD\_Au}_t$ is the gold-dollar gap
$P^{\text{mkt}}(\text{Au}) / P^{\text{FED}}(\text{Au})$ capturing
dollar-depreciation confounding, and $D^{\text{supply}}_t$ includes supply-side
dummies for the El Teniente strike (1967) and copper nationalisation (1971). The
baseline lag order is $p = 1$, following \citet{Jorda2005}'s parsimony
recommendation for small samples; $p = 2$ is estimated as a robustness check.

Horizons are set at $h = 0, 1, \ldots, 15$ for the full sample (1925--2010) and
$h = 0, \ldots, 10$ for event-window specifications centred on the Korean War
(1929--1956) and Vietnam War (1929--1971).

Two identification assumptions govern Equation~\eqref{eq:link1}:

\begin{assumption}[Exogeneity of $G^m$]
\label{ass:A1}
U.S.\ military expenditure is determined by Cold War geopolitics and is
orthogonal to Chilean macroeconomic conditions. Chile is a small open economy
that does not influence U.S.\ defence appropriations. Formally:
$\mathbb{E}[\varepsilon_{t+h} \mid \Delta G^m_t, \mathbf{X}_t] = 0$.
\end{assumption}

\begin{assumption}[Mediator exclusion]
\label{ass:A2}
U.S.\ capacity utilisation $\mu_{\text{US},t}$ is excluded from
Equation~\eqref{eq:link1}. It lies on the causal pathway
$G^m \to \mu_{\text{US}} \to P_{\text{Cu}} \to \text{ToT}$; conditioning on it
blocks the identified channel. This is a \textit{bad-control problem} in the
sense of \citet{AngristPischke2009}: including the mediator absorbs the
variation that identifies the reduced-form effect. The Military Keynesian
channel ($G^m \to \mu_{\text{US}}$) is confirmed separately in a Stage~A
regression that does not enter the Link~1 specification.
\end{assumption}

The remaining controls in Equation~\eqref{eq:link1} serve specific roles.
$\widetilde{\text{LevR}}_t$ enters in levels as a conditioning variable rather
than a control --- it captures the structural position of the balance of
payments, which may shift the intercept of the external transmission equation
without belonging to the $G^m \to \text{ToT}$ causal pathway. The gold-dollar
gap $\Delta\text{USD\_Au}_t$ absorbs variation in the terms of trade attributable
to dollar depreciation rather than to copper demand. The supply-side dummies
address discrete events that shifted Chile's export capacity independently of
world copper prices.

% ==============================================================================
\subsection{Link 2: Domestic Propagation}
\label{subsec:link2}

Link~2 identifies the domestic propagation of terms-of-trade shocks into Chilean
macroeconomic outcomes. The estimand is the impulse response of outcome $Y$ to
a terms-of-trade shock at horizon $h$:
%
\begin{equation}
\label{eq:link2}
\Delta Y_{t+h} = \alpha_h + \beta_h \,\Delta\text{ToT}_t
  + \gamma_h \,\widetilde{\text{LevR}}_t + \delta_h \,\Delta\text{USD\_Au}_t
  + \sum_{k=1}^{p} \rho_{h,k} \,\Delta Y_{t-k}
  + \sum_{k=1}^{p} \varphi_{h,k} \,\Delta\text{ToT}_{t-k}
  + \varepsilon_{t+h}
\end{equation}
%
where the outcome vector is $Y \in \{\mu_{\text{CHL}},\; u,\; \omega,\; r,\;
I\}$, comprising capacity utilisation, unemployment, the wage share, the profit
rate, and investment. The horizon is $h = 0, 1, \ldots, 10$ over the full sample
(1925--2010).

The \textit{primary outcome} is Chilean capacity utilisation $\mu_{\text{CHL}}$:
%
\begin{equation}
\label{eq:link2_mu}
\Delta\mu_{\text{CHL},t+h} = \alpha_h + \beta_h \,\Delta\text{ToT}_t
  + \gamma_h \,\widetilde{\text{LevR}}_t + \delta_h \,\Delta\text{USD\_Au}_t
  + \text{lags} + \varepsilon_{t+h}
\end{equation}
%
This prioritisation follows from the \citet{Shaikh2016} profit rate decomposition
$r = p \cdot B_r \cdot \mu \cdot (1 - \omega)$, where $p$ denotes relative
prices, $B_r$ capital productivity, $\mu$ capacity utilisation, and $\omega$ the
wage share. Capacity utilisation is the component through which external shocks
enter the profit rate; the remaining components ($\omega$, $B_r$, $p$) are
either estimated as separate outcomes or inherited as cross-essay assets from
Chapter~2. The profit rate $r$ and accumulation rate $g^K = \chi_t \cdot r_t$ ---
where $\chi_t = I_t / \Pi_t$ is the rate of capitalization \citep{Weisskopf1979,
BasuDas2017} --- are estimated as secondary outcomes whose interpretation
is disciplined by the decomposition.

\begin{assumption}[Terms-of-trade exogeneity]
\label{ass:A3}
Chile is a price-taker in the world copper market. At annual frequency,
domestic output does not feed back into the terms of trade. Copper constitutes
60--80\% of Chilean export revenues throughout the sample period, making
$\text{ToT}_t$ effectively a function of world copper prices determined in
global markets.
\end{assumption}

% ==============================================================================
\subsection{Link 3: State-Dependent Transmission}
\label{subsec:link3}

Link~3 is the central empirical contribution. It tests whether the
balance-of-payments state mediates the domestic transmission of external
shocks by conditioning the Link~2 local projections on the leverage ratio
$\widetilde{\text{LevR}}_t = e_t \cdot \text{IR}^{\text{USD}}_t / M_t$, as
defined in equation~\eqref{eq:levr_def} and whose law of motion is governed
by equation~\eqref{eq:levr_lom}.

The baseline state-dependent specification partitions the sample using an
indicator function on the leverage ratio:
%
\begin{equation}
\label{eq:link3}
\Delta Y_{t+h} = \mathbb{1}(\widetilde{\text{LevR}}_t > \tau) \cdot
  \bigl[\alpha_h^H + \beta_h^H \,\Delta\text{ToT}_t\bigr]
  + \mathbb{1}(\widetilde{\text{LevR}}_t \leq \tau) \cdot
  \bigl[\alpha_h^L + \beta_h^L \,\Delta\text{ToT}_t\bigr]
  + \gamma_h \,\Delta\text{USD\_Au}_t + \text{lags} + \varepsilon_{t+h}
\end{equation}
%
The superscripts $H$ and $L$ index the high- and low-leverage regimes
respectively. The null hypothesis is:
%
\begin{equation}
\label{eq:null_link3}
H_0 \colon \beta_h^H = \beta_h^L \quad \forall \; h
\end{equation}
%
Rejection implies that the balance-of-payments state mediates domestic
transmission of external shocks. The theoretical prediction is directional:
$\beta_h^H > 0$ (high leverage opens an expansionary corridor for terms-of-trade
improvements) and $\beta_h^L \approx 0$ or $\beta_h^L < 0$ (low leverage
constrains or reverses the domestic response).

\subsubsection*{The $\{\widetilde{\text{LevR}} \times u\}$ 2$\times$2 state space}

The full state space is defined by two binary indicators --- the leverage ratio
and the unemployment rate --- yielding a four-cell partition. Each cell
corresponds to a distinct historical configuration and theoretical prediction:

\begin{table}[ht]
\centering
\caption{State-space partition: $\{\widetilde{\text{LevR}} \times u\}$}
\label{tab:state_space}
\begin{tabular}{llllp{5.5cm}}
\hline
Cell & $\widetilde{\text{LevR}}$ & $u$ & Historical archetype & Theoretical prediction \\
\hline
$HL$ & High & Low & Korea/Vietnam boom & $\beta_h^{HL} > 0$, large; both
  corridors reinforce \\
$HH$ & High & High & Import compression & $\beta_h^{HH} > 0$ for $\mu$,
  $\approx 0$ for $\omega$ \\
$LL$ & Low & Low & UP 1970--73 (explosive) & Amplification via Corridor~B
  feedback; double-squeeze operative \\
$LH$ & Low & High & Post-1973 shock therapy & $\beta_h^{LH} \approx 0$;
  Corridor~B dismantled \\
\hline
\end{tabular}
\end{table}

The $LL$ cell is the theoretically critical configuration: low leverage
(external constraint binding) combined with low unemployment (reserve army
exhausted). This is the Unidad Popular configuration. The theory predicts that
terms-of-trade shocks in this cell are amplified through Corridor~B feedback
loops --- the wage share rises as the reserve army is depleted (Marx), aggregate
consumption expands \citep{Kalecki1943}, and the resulting import demand and
monetary expansion erode the leverage ratio further (the \textit{double squeeze}
through channels $P$ and $Q$ in Figure~\ref{fig:causal_graph}).

\subsubsection*{Nesting strategy and degrees of freedom}

The state-dependent specification is estimated at three levels of increasing
granularity to manage the degrees-of-freedom constraint imposed by the sample
size ($N \approx 85$ usable observations).

\textit{Level (a): LevR-only threshold} (Equation~\ref{eq:link3}). Two regimes,
estimating separate intercept-slope pairs for high and low leverage. This is the
primary specification.

\textit{Level (b): $u$-only threshold.} The labour-market state conditions the
distributional response:
%
\begin{equation}
\label{eq:u_threshold}
\Delta\omega_{t+h} = \mathbb{1}(u_t < \tau_u) \cdot
  \bigl[\alpha_h^{TL} + \beta_h^{TL} \,\Delta\text{ToT}_t\bigr]
  + \mathbb{1}(u_t \geq \tau_u) \cdot
  \bigl[\alpha_h^{SL} + \beta_h^{SL} \,\Delta\text{ToT}_t\bigr]
  + \gamma_h \,\widetilde{\text{LevR}}_t + \text{lags} + \varepsilon_{t+h}
\end{equation}
%
with null $H_0 \colon \beta_h^{TL} = \beta_h^{SL}$ (no reserve-army state
dependence). The superscripts $TL$ and $SL$ denote tight and slack labour
markets. The threshold $\tau_u$ is calibrated from the historical record at
approximately 4.5--5.0\%, with the UP exhaustion floor at 3.5--4.0\%.

\textit{Level (c): Joint 4-cell.} The full interaction:
%
\begin{equation}
\label{eq:joint_4cell}
\Delta\omega_{t+h} = \sum_{s \in \{HL, HH, LL, LH\}}
  \mathbb{1}(s_t = s) \cdot
  \bigl[\alpha_h^s + \beta_h^s \,\Delta\text{ToT}_t\bigr]
  + \text{controls} + \varepsilon_{t+h}
\end{equation}
%
With four cell-specific intercept-slope pairs, the model requires estimation of
8 regime parameters plus controls per horizon. At $p = 1$ with 2 controls, this
yields approximately 12 parameters per horizon against 85 observations ---
a degrees-of-freedom ratio of roughly 7:1. This is feasible but tight. The joint
4-cell specification is therefore estimated as a robustness exercise; the primary
identification rests on the nested single-threshold models
(Equations~\ref{eq:link3} and \ref{eq:u_threshold}).

Two additional assumptions govern the state-dependent design:

\begin{assumption}[LevR as state variable]
\label{ass:A4}
The leverage ratio conditions the intensity of external-shock transmission but
is not a policy instrument chosen endogenously in response to contemporaneous
shocks. It is a slow-moving stock variable reflecting the accumulated position
of the balance of payments. Its components --- exchange rate, international
reserves, and the monetary base --- adjust at frequencies slower than the
annual shocks being identified.
\end{assumption}

\begin{assumption}[ECT$_m$ regime identification]
\label{ass:A5}
The error-correction term $\text{ECT}_{m,t}$ from the Chapter~2 threshold
vector error-correction model provides an independent regime-identification
device. When $\text{ECT}_{m,t} > \hat{\lambda}$, the balance-of-payments
constraint is binding. This cross-essay asset complements
$\widetilde{\text{LevR}}_t$ by distinguishing structural position (the level
of the leverage ratio) from binding episodes (threshold exceedance in the
error-correction dynamics).\footnote{%
  The conservative default $\hat{\lambda} = 0$ is used in this draft, pending
  calibration from the Stage~A Chile estimation. Decision locked 2026-04-08.}
\end{assumption}

\subsubsection*{Threshold calibration}

The $\widetilde{\text{LevR}}$ threshold $\tau_{\text{LevR}}$ is computed as the
within-regime median, calculated separately for the pre- and post-1971
subsamples. This approach avoids look-ahead bias: the threshold uses only
observations available within each structural regime defined by the Bretton
Woods breakdown. The ECT$_m$ binary provides a complementary identification
device that captures whether the constraint is actively binding in a given year,
irrespective of the structural level of $\widetilde{\text{LevR}}$.

% ==============================================================================
\subsection{Identification and Inference}
\label{subsec:identification}

The identification strategy rests on six assumptions, collected here for
reference.

Assumptions~\ref{ass:A1}--\ref{ass:A3} govern the sequential chain: $G^m$
exogeneity (A1) ensures that U.S.\ military spending is a valid instrument for
external shocks; mediator exclusion (A2) prevents the bad-control problem from
blocking the identified channel; terms-of-trade exogeneity (A3) ensures that the
domestic propagation equations are not contaminated by reverse causality from
Chilean output to world copper prices. Assumptions~\ref{ass:A4}--\ref{ass:A5}
govern the state-dependent design: $\widetilde{\text{LevR}}$ as a predetermined
state variable (A4) and the ECT$_m$ regime-identification device (A5) jointly
ensure that the conditioning partition is not endogenous to the shocks being
estimated.

The sixth assumption provides the structural-break falsification test:

\begin{assumption}[Post-1973 placebo]
\label{ass:A6}
The Pinochet regime dismantled the institutional mechanisms constituting
Corridor~B: unions were suppressed, wages repressed, and the economy was
reintegrated into international capital markets under floating exchange rates
and liberalised trade. If the Link~2 coefficients are significant in the
pre-1973 sample but insignificant in the post-1973 sample, while Link~1
persists across both periods, the domestic propagation mechanism is confirmed
as historically contingent rather than a statistical artefact of the full
sample. A secondary sub-period break at 1982 (the debt crisis) tests the
stability of the post-1973 null.
\end{assumption}

\subsubsection*{Inference}

Every impulse response function is reported with \textit{dual-band} 95\%
confidence intervals. The inner (dark) band uses HC3 heteroskedasticity-consistent
standard errors, which are robust to conditional heteroskedasticity without
imposing a serial-correlation structure and incorporate a conservative
small-sample correction. The outer (light) band uses Newey-West HAC standard
errors with bandwidth $h + 1$, accounting for the moving-average structure
inherent in overlapping local-projection residuals at horizon $h$, following
\citet{Jorda2005}.

Both bands are needed because local projections at horizon $h$ generate
mechanically correlated residuals across adjacent horizons: $Y_{t+h}$ and
$Y_{t+h+1}$ share $h$ observations in the dependent variable. HC3 is valid
per-horizon but ignores this cross-horizon dependence; NW/HAC corrects for it
but may over-reject in small samples. Inferences that hold under both bands are
treated as reliable; those that hold only under HC3 are flagged as provisional.

A third robustness layer is provided by the \textit{wild bootstrap}:
%
\begin{equation}
\label{eq:wild_bootstrap}
\text{95\% CI}^{\text{WB}} \colon \quad
  1{,}000 \text{ bootstrap replications, Rademacher weights}
\end{equation}
%
The wild bootstrap provides inference robust to both heteroskedasticity and the
non-standard error structure of LP residuals under small $N$. It is deployed
when HC3 and NW bands diverge, serving as the tie-breaking criterion.

\subsubsection*{The $M(s)$ multiplier}

The compound causal intensity multiplier $M(s)$ aggregates chain-specific
impulse responses, weighted by their temporal overlap, at state $s$:
%
\begin{equation}
\label{eq:multiplier}
M(s) = \sum_{(l, l') \in \mathcal{P}^*}
  \tilde{\beta}_l(s) \cdot \tilde{\beta}_{l'}(s) \cdot |O_{l,l'}|
\end{equation}
%
where $\mathcal{P}^* = \{(D, E),\; (D, F),\; (D, I)\}$ is the set of primary
chain pairs, $\tilde{\beta}_l(s) = \beta_l(s) / \sigma_l$ is the
SD-normalised impulse response for chain $l$ at state $s$, and $O_{l,l'} =
C_l \cap C_{l'}$ is the calendar-year overlap between the active sets of chains
$l$ and $l'$. $M(s)$ is a \textit{results-section diagnostic}: it is computed ex
post from estimated impulse response functions and is not a pre-specified test
statistic. Its role is interpretive --- it quantifies whether the Vietnam and UP
episodes exhibit compound causation (large $|O_{l,l'}|$, both $\tilde{\beta}$
significant) or merely sequential causation (non-overlapping active windows).
This operationalises Robinson's (1980) Historical Time principle: causality
requires temporal co-activation, not merely statistical association.

% ==============================================================================
\subsection{Exercises and Robustness}
\label{subsec:exercises}

The estimation programme is organised into two exercises and a set of robustness
checks.

\subsubsection*{Exercise A: Full-sample aggregate (1925--2010)}

Exercise~A estimates the five aggregate outcomes under Equations~\eqref{eq:link1}
and \eqref{eq:link2}: (i)~$\Delta\text{ToT}$ (Link~1, Chain~D);
(ii)~$\Delta\mu_{\text{CHL}}$ (Link~2, primary outcome);
(iii)~$\Delta\omega$ (distributional, Chain~I);
(iv)~$\Delta r$ (profitability, Chain~E);
(v)~$\Delta I/K$ (accumulation, Chain~E). The full sample pools across all
structural regimes; state-dependent conditioning is applied in the Link~3
specifications.

\subsubsection*{Exercise B: Sectoral EOD outcomes (1957--1979)}

Exercise~B estimates sectoral employment-output-distribution outcomes
$\{u_s, \omega_s\}$ under Chain~E$'$ for the period spanning the Lewis turning
point through the UP crisis. The restricted sample isolates the window in which
the reserve army is actively depleting \citep{Lewis1954}, the \citet{Goodwin1967}
distributive cycle is operative, and the political dimension of full employment
\citep{Kalecki1943} becomes binding. This exercise tests whether the aggregate
reserve-army mechanism is driven by specific sectors or operates as a
macroeconomic regularity.

\subsubsection*{Robustness checks}

Five robustness checks are conducted:

\begin{itemize}
  \item \textit{$\widetilde{\text{LevR}}_{\text{gold}}$.} A gold-equivalent
    leverage ratio addressing the 1971 Bretton Woods regime break. This variant
    ensures that the state-dependent results are not driven by the shift from
    fixed to floating exchange rates.

  \item \textit{Event-window vs.\ pooled.} Korea-centred (1929--1956) and
    Vietnam-centred (1929--1971) event windows are estimated alongside the
    full-sample pooled specification. Divergence between event-window and pooled
    estimates indicates structural instability.

  \item \textit{Lag-length sensitivity.} $p = 2$ is estimated against the
    $p = 1$ baseline across all links.

  \item \textit{NW bandwidth sensitivity.} Alternative HAC bandwidths
    $\{h,\; h+1,\; h+2\}$ are compared to assess sensitivity of the outer
    confidence band.

  \item \textit{Sub-period structural break.} Pre/post-1973 sample splits
    implement Assumption~\ref{ass:A6}. A secondary break at 1982 tests whether
    the post-1973 null is stable or disrupted by the debt crisis.
\end{itemize}

% ==============================================================================
% End of Section 4 — Empirical Strategy
% ==============================================================================


%</ch3body>

\pagebreak
% REF References
\label{ch3:references}
\bibliographystyle{apalike}
\bibliography{../references}


\pagebreak
% In Chapter3_PEUP.tex, before the \input:
\appendix
%<*ch3appendix>
% ============================================================
% APPENDIX: Balance-of-Payments Decomposition and Leverage
%           Ratio Algebra
% Cross-referenced from Section C (eq. levr_lom, eq. binding_condition)
% Contains: Steps 1–6 of the formalization
% ============================================================


\section{The Leverage Ratio Embedded in the Balance-of-Payments Constraint}
\label{app:bop_levr}

This appendix develops the full accounting decomposition underlying the
leverage ratio's law of motion (equation~\ref{eq:levr_lom}) and the
constraint-binding condition (equation~\ref{eq:binding_condition})
presented in Section~\ref{sec:C2_bop_dynamics}. The derivation proceeds
in six steps: definition of the leverage ratio, decomposition of the
current and capital accounts, the general balance-of-payments identity,
substitution into the leverage ratio, and a reading of the constraint
under alternative scenarios.

\subsection{The leverage ratio in nominal terms}
\label{app:levr_nominal}

Let all variables be denominated in current domestic currency (CLP) at
time $t$. The exchange rate convention is $e_t$ = CLP per USD, so that
depreciation implies $e_t \uparrow$.

Define:
\begin{itemize}
\item $IR_t$: stock of international reserves, denominated in USD
\item $e_t$: nominal exchange rate (CLP/USD)
\item $M_t$: currency in circulation---the central bank's active domestic
  monetary liability (CLP)
\end{itemize}

The central bank leverage ratio in nominal terms:
%
\begin{equation}
\label{app:eq:levr_def}
\widetilde{\text{LevR}}_t \;\equiv\; \frac{e_t \cdot IR_t}{M_t}
\end{equation}
%
This is a dimensionless ratio: the peso value of the central bank's
world-money holdings relative to its domestic monetary liability issuance.
In growth rates:
%
\begin{equation}
\label{app:eq:levr_growth}
g_{\widetilde{\text{LevR}},t} \;=\; g_{e,t} \;+\; g_{IR,t} \;-\; g_{M,t}
\end{equation}


\subsection{Current account decomposition}
\label{app:current_account}

The current account in domestic currency:
%
\begin{equation}
\label{app:eq:CA}
CA_t = e_t \cdot P_{M,t}^{USD}\!\left(\text{ToT}_t \cdot X_t
       - M_t^{goods}\right)
       - e_t \cdot r_t \cdot D_t^{ext} + T_t
\end{equation}
%
where $\text{ToT}_t \equiv P_{X,t}^{USD} / P_{M,t}^{USD}$ is the net
barter terms of trade. The three components are:

\begin{table}[h]
\centering
\small
\caption{Components of the Current Account}
\begin{tabular}{lll}
\toprule
Term & Content & Character \\
\midrule
$e_t P_{M,t}^{USD}(\text{ToT}_t \cdot X_t - M_t^{goods})$
  & Trade balance at import prices
  & ToT exogenous for small open economy \\
$e_t \cdot r_t \cdot D_t^{ext}$
  & External debt service
  & Largely predetermined \\
$T_t$
  & Net current transfers
  & Exogenous \\
\bottomrule
\end{tabular}
\label{app:tab:CA}
\end{table}

The trade balance is expressed at import prices so that the terms-of-trade
effect enters multiplicatively through $\text{ToT}_t$. A terms-of-trade
improvement ($\text{ToT}_t \uparrow$) raises the trade balance for given
export and import volumes. Debt service is largely predetermined by the
existing stock of external obligations $D_t^{ext}$ and the interest rate
$r_t$ on external debt. Net transfers $T_t$ are exogenous.


\subsection{Capital account decomposition}
\label{app:capital_account}

The capital account in domestic currency:
%
\begin{equation}
\label{app:eq:KA}
KA_t = e_t \cdot \Delta D_t^{ext} + e_t \cdot FDI_t + e_t \cdot \Pi_t
\end{equation}
%
where $\Delta D_t^{ext}$ is net external borrowing, $FDI_t$ is net foreign
direct investment, and $\Pi_t$ denotes net portfolio and short-term capital
flows, with $\Pi_t < 0$ representing capital flight. The capital account
is the channel through which the profit squeeze translates into reserve
drainage: when the domestic profit rate falls and the external constraint
is simultaneously tightening, domestic capital has both the motive (low
profitability) and the means (still-available foreign exchange) to move
resources abroad, compounding the leverage deterioration from the current
account side.


\subsection{The general balance-of-payments identity}
\label{app:bop_identity}

The change in the peso value of international reserves equals the sum of
the current and capital account balances:
%
\begin{equation}
\label{app:eq:bop}
\Delta(e_t \cdot IR_t) = CA_t + KA_t \equiv \text{BoP}_t
\end{equation}
%
Three structural channels operate through this identity:

\begin{enumerate}
\item \textbf{Trade channel:} Terms-of-trade deterioration compresses
  export revenue relative to imports, draining reserves through the
  current account.
\item \textbf{Financial channel:} Debt service outflows or capital flight
  drain reserves independently of trade performance.
\item \textbf{Valuation channel:} Exchange rate depreciation raises the
  CLP value of existing USD reserves---a pure balance sheet effect that
  changes $\widetilde{\text{LevR}}_t$ without altering the underlying
  real position.
\end{enumerate}


\subsection{Substitution into the leverage ratio}
\label{app:levr_substitution}

From the definition $e_t \cdot IR_t = \widetilde{\text{LevR}}_t \cdot
M_t$, substituting into equation~\eqref{app:eq:bop}:
%
\begin{equation}
\label{app:eq:levr_sub}
\Delta\!\left(\widetilde{\text{LevR}}_t \cdot M_t\right) = CA_t + KA_t
\end{equation}
%
Expanding the left-hand side using the product rule and rearranging:
%
\begin{equation}
\label{app:eq:levr_lom}
\boxed{\;
\Delta\widetilde{\text{LevR}}_t \;=\;
\frac{CA_t + KA_t}{M_t}
\;-\; \widetilde{\text{LevR}}_t \cdot g_{M,t}
\;}
\end{equation}
%
This is equation~\eqref{eq:levr_lom} in the main text. The first term is
the net balance-of-payments flow per unit of domestic monetary liability.
The second term---$\widetilde{\text{LevR}}_t \cdot g_{M,t}$---is the
monetary expansion multiplier: the existing leverage ratio scaled by the
growth rate of domestic money. The product structure implies that the same
rate of monetary expansion ($g_{M,t}$) produces a larger absolute drain on
the leverage position when $\widetilde{\text{LevR}}_t$ is high, and a
smaller absolute drain when it is low---but in the latter case the margin
of safety is also smaller, so that even modest monetary expansion can
push the ratio below the binding threshold.


\subsection{Reading the constraint: cases and scenarios}
\label{app:constraint_cases}

The leverage position is deteriorating whenever:
%
\begin{equation}
\label{app:eq:binding}
\Delta\widetilde{\text{LevR}}_t < 0 \;\iff\;
\frac{CA_t + KA_t}{M_t} < \widetilde{\text{LevR}}_t \cdot g_{M,t}
\end{equation}
%
This is equation~\eqref{eq:binding_condition} in the main text. A falling
leverage ratio does not mean the constraint is binding---it means the
constraint is tightening, the distance to the binding threshold is
shrinking. Whether the constraint actually binds depends on the
\textit{level} of $\widetilde{\text{LevR}}_t$.

\paragraph{Case 1: Balance-of-payments deficit}
($CA_t + KA_t < 0$). The constraint tightens trivially. Two forces
compound: reserves are depleted by the BoP deficit (numerator drains),
and monetary expansion erodes the coverage ratio independently
(denominator dilutes).

\paragraph{Case 2: Balance-of-payments surplus}
($CA_t + KA_t > 0$). The constraint tightens whenever the surplus is
insufficient to cover the monetary expansion multiplier:

\begin{table}[h]
\centering
\small
\begin{tabular}{ll}
\toprule
Condition & $\Delta\widetilde{\text{LevR}}_t$ \\
\midrule
Surplus $<$ multiplier threshold & $< 0$ --- tightening despite surplus \\
Surplus $=$ multiplier threshold & $= 0$ --- stable \\
Surplus $>$ multiplier threshold & $> 0$ --- loosening \\
\bottomrule
\end{tabular}
\caption{Leverage dynamics under balance-of-payments surplus.}
\label{app:tab:surplus}
\end{table}

\paragraph{Historical scenarios.}
Three configurations illustrate the constraint's operation across the
regimes that the local projections specification in
Section~\ref{sec:C2_bop_dynamics} conditions on:

\begin{table}[h]
\centering
\small
\begin{tabular}{lccll}
\toprule
Scenario & $e_t \cdot IR_t$ & $M_t$ &
  $\Delta\widetilde{\text{LevR}}_t$ & Mechanism \\
\midrule
A --- Fiscal monetisation
  & Stable & Rising & Falling & Inside pressure \\
B --- UP 1970--73
  & Falling & Rising fast & Falling fast & Double compression \\
C --- Shock therapy post-1973
  & Slowly rising & Falling & Rising & Denominator destruction \\
\bottomrule
\end{tabular}
\caption{Leverage ratio dynamics under three historical configurations.
Scenario~A: the constraint tightens from the liability side alone.
Scenario~B: reserve drain and monetary expansion compound.
Scenario~C: leverage recovers through denominator compression, not
reserve accumulation.}
\label{app:tab:scenarios}
\end{table}

Scenario~A demonstrates the inside-pressure channel in isolation: even
with a stable or positive external position, fiscal monetisation expands
$M_t$ and the leverage ratio falls. Scenario~B is the double compression
that characterises the Unidad Popular period: high $g_{M,t}$ from
wage-fiscal expansion compounds with reserve drainage through the current
account, so that the constraint became acute through the multiplier before
the full reserve crisis materialised. Scenario~C shows the reverse: the
post-1973 austerity programme restored the leverage ratio not through
reserve accumulation but through denominator compression---the contraction
of domestic monetary liabilities.

The multiplier's asymmetric structure across leverage states is summarised
in Table~\ref{app:tab:multiplier}.

\begin{table}[h]
\centering
\small
\begin{tabular}{lll}
\toprule
LevR level & Multiplier
  $\widetilde{\text{LevR}}_t \cdot g_{M,t}$ & Consequence \\
\midrule
High   & Large    & Large absolute drain---from safe position \\
Medium & Moderate & Moving toward constraint \\
Low    & Small    & Small absolute drain---potentially decisive
                    near threshold \\
\bottomrule
\end{tabular}
\caption{Multiplier structure across leverage states.}
\label{app:tab:multiplier}
\end{table}

The state-dependent local projections specification in
Section~\ref{sec:C2_bop_dynamics} exploits precisely this asymmetry:
identical terms-of-trade shocks are expected to produce different domestic
impulse responses depending on the leverage state at the moment the shock
arrives.

%</ch3appendix>

\end{document}
